\documentclass[12pt,a4paper]{article}
\usepackage[margin=2.4cm]{geometry}
\usepackage{amsmath,amssymb,bm,mathtools}
\usepackage{fontspec}
\setmainfont{TeX Gyre Termes}
\setsansfont{TeX Gyre Heros}
\setmonofont{Latin Modern Mono}[Scale=0.9]
\usepackage{graphicx,float,caption,subcaption,booktabs,longtable,array,multirow,tabularx}
\usepackage{xcolor}
\usepackage{setspace}
\usepackage{lineno}
\usepackage{microtype}
\usepackage{xurl}
\usepackage{siunitx}
\usepackage{enumitem}
\usepackage{fancyhdr}
\usepackage{titlesec}
\titleformat{\paragraph}[runin]{\normalfont\itshape}{}{0pt}{}
\titlespacing*{\paragraph}{0pt}{1.1ex plus .2ex}{0.7em}
\usepackage[authoryear,round,semicolon,sort&compress]{natbib}
\definecolor{lk}{RGB}{0,70,140}
\usepackage[colorlinks=true,linkcolor=lk,citecolor=lk,urlcolor=lk,breaklinks=true,pdfborder={0 0 0}]{hyperref}
\usepackage[nameinlink,capitalise]{cleveref}
\sisetup{detect-all,separate-uncertainty=false}
\setstretch{1.5}
\setlength{\parskip}{3pt}
\captionsetup{font=small,labelfont=bf,justification=centering,singlelinecheck=true}
\titleformat{\section}{\large\bfseries}{\thesection.}{0.5em}{}
\titleformat{\subsection}{\normalsize\bfseries}{\thesubsection}{0.5em}{}
\titleformat{\subsubsection}{\normalsize\itshape}{\thesubsubsection}{0.5em}{}
\pagestyle{fancy}\fancyhf{}\fancyfoot[C]{\small\thepage}\renewcommand{\headrulewidth}{0pt}
\setlength{\headheight}{14pt}
\bibliographystyle{plainnathl}
\newcommand{\doihl}[1]{doi: \href{https://doi.org/#1}{\nolinkurl{#1}}}
\newcommand{\dd}{\mathrm{d}}
\newcommand{\ET}{\mathit{ET}}
\newcommand{\Ent}{\mathit{ET}_{\mathrm{nt}}}
\newcommand{\Ks}{K_{\mathrm{s}}}
\newcommand{\Se}{S_{\mathrm{e}}}
\newcommand{\TAW}{\mathit{TAW}}
\newcommand{\repourl}{https://github.com/halilunaziru73-creator/Coupled-Water-Fluxes-Pore-Scale-Dynamics-Precision-Irrigation-Mulching-Feedback-Loops}
\newcommand{\repo}{\href{\repourl}{repository}}
\newcommand{\dataurl}{\repourl/tree/main/Data}
\newcommand{\codeurl}{\repourl/tree/main/Manuscripts}
\newcommand{\softurl}{\repourl/tree/main/Software}
\newcommand{\datalink}{\href{\dataurl}{Data}}
\newcommand{\codelink}{\href{\codeurl}{Code}}
\newcommand{\softlink}{\href{\softurl}{Software}}
\newcommand{\manuallink}{\href{\manualurl}{Manual}}
\newcommand{\downloadlink}{\href{\exeurl}{SWAFL.exe}}
\newcommand{\figdir}{../figures}
\newcommand{\manualurl}{https://github.com/halilunaziru73-creator/Coupled-Water-Fluxes-Pore-Scale-Dynamics-Precision-Irrigation-Mulching-Feedback-Loops/raw/main/Software/v1.1.0/SWAFL_Desktop/docs/manual_v1.1.0/SWAFL_Manual_v1.1.0.pdf}
\newcommand{\exeurl}{https://github.com/halilunaziru73-creator/Coupled-Water-Fluxes-Pore-Scale-Dynamics-Precision-Irrigation-Mulching-Feedback-Loops/releases/latest/download/SWAFL.exe}
\newcommand{\authorblock}{%
\begin{center}
{\large\textbf{Naziru Halilu}}$^{1,2,*}$\\[3pt]
{\small $^{1}$ Agricultural University of Athens, Athens, Greece\\
$^{2}$ Public University of Navarre, Pamplona, Spain\\
ORCID: \href{https://orcid.org/0009-0006-2324-9599}{0009-0006-2324-9599}\\
$^{*}$ Correspondence: \href{mailto:halilunaziru73@gmail.com}{halilunaziru73@gmail.com}}
\end{center}}
\newcommand{\csurl}{https://github.com/halilunaziru73-creator/Coupled-Water-Fluxes-Pore-Scale-Dynamics-Precision-Irrigation-Mulching-Feedback-Loops/tree/main/Data/Comprehensive_Summary}
\newcommand{\cslink}{\href{\csurl}{\nolinkurl{Data/Comprehensive_Summary}}}
\newcommand{\csavail}{\par\smallskip\textbf{Full simulation outputs.} Every value of every log of every run (the annual water and nitrogen balances, the harvests and the crop state, year by year), one compressed table per experiment (\texttt{E00\_Baselines} to \texttt{E16\_Edge\_Cases}), the one-row-per-run summary of all 10\,190 runs (\texttt{00\_summary\_all\_simulations}), the per-experiment summary (\texttt{00\_SUMMARY\_ALL\_EXPERIMENTS}) and the script that made them (\texttt{summarize\_results3.py}) are in the folder \cslink{} of the repository (\url{\csurl}); their 20-year means reproduce the run table within 2\,\%.}
% Data and code availability of the two analysis papers (P1, P2)
\newcommand{\dataavail}{%
\textbf{Data.} The complete Daisy ensemble outputs (all 10\,190 runs, raw and derived) are archived on Zenodo, \doihl{10.5281/zenodo.23043626}. The raw run table used here, and a linked list of every data file used, are in the \datalink{} folder of the project repository (\url{\dataurl}).\csavail
}
% Software, data and code availability of the software paper (P3)
\newcommand{\softavail}{%
\textbf{Software.} SWAFL Desktop (version 1.1.0) is open source; the \softlink{} source is in the project repository (\url{\softurl}) and the program for Windows is \downloadlink{} (\url{\exeurl}).
\par\smallskip\textbf{Manual.} The \manuallink{} (\url{\manualurl}) documents every control, view and equation and labels the provenance of all 134 numbered equations; it can be reviewed or downloaded as a PDF.
\par\smallskip\textbf{Data.} The complete Daisy ensemble outputs (all 10\,190 runs, raw and derived) are archived on Zenodo, \doihl{10.5281/zenodo.23043626}. The raw run table, and a linked list of every data file used, are in the \datalink{} folder of the repository (\url{\dataurl}).\csavail
}

\usepackage{pdflscape}
\setlength{\emergencystretch}{3em}

% A shared availability statement
\newcommand{\suiteavail}{%
\textbf{Software.} SWAFL Desktop (version 1.1.0) is open source; the \softlink{} source is in the project repository (\url{\softurl}) and the program for Windows is \downloadlink{} (\url{\exeurl}). The whole workbench also runs in a browser, with nothing to install (link in the repository README).
\par\smallskip\textbf{Manual.} The \manuallink{} (\url{\manualurl}) documents every control, screen and equation of the release, with labelled screenshots.
\par\smallskip\textbf{Data.} The complete Daisy ensemble outputs (all 10\,190 runs, raw and derived) are archived on Zenodo, \doihl{10.5281/zenodo.23043626}; the run table and a summary of every experiment are in the \datalink{} folder (\url{\dataurl}).\csavail
}

\renewcommand{\thesection}{S\arabic{section}}
\renewcommand{\theequation}{S\arabic{equation}}
\renewcommand{\thetable}{S\arabic{table}}
\renewcommand{\thefigure}{S\arabic{figure}}
\hypersetup{pdftitle={Supplementary Material: SWAFL Desktop as an Offline and Equation-Traceable Decision-Support Suite with a Graphical Workbench and a Verified Bridge for the Daisy Model},pdfauthor={Naziru Halilu}}
\begin{document}
\begin{center}{\large\bfseries Supplementary Material\par}\vspace{4pt}
{\Large\bfseries SWAFL Desktop as an Offline and Equation-Traceable Decision-Support Suite with a Graphical Workbench and a Verified Bridge for the Daisy Model\par}\vspace{6pt}
{\normalsize Derivations of the equations, the mechanisms behind the results, the data, and the origin of every table and figure\par}\end{center}
\authorblock
\tableofcontents
\clearpage

\section{Derivations of the equations}
\subsection{Normalised spectral indices}
Green leaves absorb red light for photosynthesis and scatter near-infrared light in the leaf mesophyll. The normalised difference $\mathrm{NDVI}=(\rho_N-\rho_R)/(\rho_N+\rho_R)$ is a contrast of the two that is bounded in $[-1,1]$ and cancels a common multiplicative illumination factor $k$: $(k\rho_N-k\rho_R)/(k\rho_N+k\rho_R)$ equals NDVI. SAVI adds a soil term $L$ to the denominator to reduce the influence of bright soil, $(1+L)(\rho_N-\rho_R)/(\rho_N+\rho_R+L)$; EVI corrects for the atmosphere with the blue band and for the canopy background; the water indices use the shortwave-infrared band, which liquid water absorbs.

\subsection{Crop coefficient through the season}
The FAO-56 curve is piecewise linear in time: constant $K_{ini}$ during the initial stage $L_1$, a linear rise to $K_{mid}$ over the development stage $L_2$, constant over the mid-season $L_3$ and a linear fall to $K_{end}$ over the late stage $L_4$ \citep{Allen1998}. The crop evapotranspiration is $\ET_c=K_c\,\ET_0$.

\subsection{Daily root-zone balance of each map cell}
For each cell, the root-zone depletion follows $D_{r,i}=D_{r,i-1}-(P-\mathrm{RO})_i-I_i+\ET_{c,i}+\mathrm{DP}_i$, bounded by $0\le D_{r}\le\TAW$; deep percolation occurs only when the depletion would fall below zero. Irrigation is required when $D_r$ exceeds the readily available water $p\,\TAW$; the net depth to apply is $D_r$ \citep{Allen1998}.

\subsection{The Richards equation of the soil column}
Mass conservation of water in a layer of thickness $\dd z$ gives $\partial\theta/\partial t=-\partial q/\partial z-S$, with $S$ the root uptake. The Darcy--Buckingham flux with $z$ positive downward and $h$ the suction head is $q=K(h)(1-\partial h/\partial z)$: gravity drives water down and a suction gradient drives it towards the drier soil. With $\theta(h)$ and $K(h)$ from the van~Genuchten--Mualem functions the equation is closed and is solved implicitly with Picard iteration in the column of the workbench.

\subsection{Nitrogen and chemical budgets}
For each season, the change of the mineral nitrogen store equals inputs minus outputs: fertiliser, deposition and mineralisation in; uptake, leaching, denitrification and volatilisation out. For a chemical, $\Delta M=A-U-L-R-X$ (applied, taken up, leached by every route, run off, transformed); the bridge reads each term from the Daisy logs and the closure of this budget is its test.

\subsection{Splitting denitrification into N$_2$O and N$_2$}
Let $r$ be the ratio of N$_2$ to N$_2$O in the gaseous product of denitrification $D$. Then $D_{\mathrm{N_2O}}+D_{\mathrm{N_2}}=D$ and $D_{\mathrm{N_2}}=r\,D_{\mathrm{N_2O}}$, so $D_{\mathrm{N_2O}}(1+r)=D$, giving $D_{\mathrm{N_2O}}=D/(1+r)$ and $D_{\mathrm{N_2}}=rD/(1+r)=D/(1+1/r)$, the two expressions of the paper. The ratio is the smaller of a nitrate and a carbon-dioxide limited ratio, scaled by a water factor, so wetter, more active soils give relatively more N$_2$.

\subsection{Terrain, flow and wetness}
On a filled elevation grid each cell drains to the steepest of its eight neighbours (D8). The upslope area $A$ of a cell is the sum of the areas that drain through it, and the specific area is $a=A/w$ with $w$ the cell width. The topographic wetness index $\ln(a/\tan\beta)$ is large where much water arrives and the slope $\beta$ is small, which marks where the soil stays wet.

\subsection{Soil loss}
The revised universal soil loss equation multiplies the rain erosivity $R$, the soil erodibility $K$, the slope length and steepness $LS$, the cover $C$ and the practice $P$: $A=R\,K\,LS\,C\,P$ in t\,ha$^{-1}$\,yr$^{-1}$ \citep{Renard1994}.

\subsection{Run budget of an experiment}
For $k$ factors with $n_1,\dots,n_k$ levels the full design has $\prod n_i$ runs. When this exceeds the budget $B$ set by the user, a Latin-hypercube choice of $B$ runs is made: each factor's levels are used as evenly as possible ($B/n_i$ times each), so every level of every factor appears and the main effects remain estimable while the number of runs never exceeds $B$.

\section{Why the results happen}
\paragraph{Why every number can be traced.} Each value shown is computed by a function that carries the number of its equation and a label of origin (literature, fitted to data, engineered); the interface shows both, so the origin is a property of the code, not of the documentation.
\paragraph{Why a Daisy layer draws nothing before a Daisy run.} The rule of origin forbids a layer to show a Daisy result from any other source; without a run there is no result, so the layer says so rather than showing an example.
\paragraph{Why the compact designs lose some shape of a curve.} A design with fewer levels samples a curve at fewer points; the marginal curve is then drawn through those points and misses curvature between them, which is the loss measured.
\paragraph{Why the soil column is slow for a ponded fine soil.} With ponding the top boundary is a fixed head and the fine soil has a very small conductivity and a steep head gradient at the surface; the implicit solver then needs small time steps and many iterations to converge.
\paragraph{Why the denitrification defect was found.} The check compared the shared function with the stored Daisy results over 792 cases; one case class disagreed, which pointed to a wrong branch in the function, corrected in the release.
\paragraph{Why runs on a horizon of the Daisy samples failed.} The sample file also defined logs that referred to a model of its original run; Daisy stopped on the unknown name. Taking only the definitions a run needs removed the error.
\paragraph{Why the routes from Daisy to a decision differ in time by orders of magnitude.} Running Daisy solves the Richards equation and a crop for every day of every run; the meta-model looks up and combines stored outcomes, which takes microseconds.

\section{The data}
\Cref{tab:data} lists the 17 experiments of the data folder and \cref{fig:data} shows their outputs. The tables that follow are the result files of the analyses of the paper, printed in full or up to 220 rows; each is also in the workbook \nolinkurl{Data/Manuscript_Data/SWAFL_Manuscripts_Data.xlsx}, and every run is in \cslink.
\begin{table}[H]\centering\scriptsize
\caption{The data of this study: the 17 Daisy experiments of \texttt{Data/Comprehensive\_Summary} (10190 runs, 10041 valid; 20 years, 1980--1999, of Taastrup weather). Means of the valid runs; $Y$ yield (Mg\,ha$^{-1}$), $I$ irrigation, $D$ drainage at 1\,m and $\ET$ (mm\,yr$^{-1}$), $N_L$ nitrate leached (kg\,N\,ha$^{-1}$\,yr$^{-1}$), $t$ run time of Daisy (s). Every row is one file of the folder, linked.}
\label{tab:data}
\setlength{\tabcolsep}{3pt}
\begin{tabular}{@{}llrrp{3.5cm}rrrrrr@{}}
\toprule
Exp. & File & Runs & Valid & Factors varied & $Y$ & $I$ & $D$ & $\ET$ & $N_L$ & $t$\\\midrule
E0 & \href{\repourl/blob/main/Data/Comprehensive_Summary/E00_Baselines.csv.gz}{\texttt{Baselines}} & 10 & 10 & soil & 2.19 & 0 & 55 & 418 & 158 & 231\\
E1 & \href{\repourl/blob/main/Data/Comprehensive_Summary/E01_Doses.csv.gz}{\texttt{Doses}} & 200 & 189 & soil, rate, dose & 5.13 & 198 & 98 & 564 & 79 & 195\\
E2 & \href{\repourl/blob/main/Data/Comprehensive_Summary/E02_TriggerPressure.csv.gz}{\texttt{TriggerPressure}} & 200 & 197 & soil, trigger & 4.70 & 169 & 80 & 559 & 90 & 139\\
E3 & \href{\repourl/blob/main/Data/Comprehensive_Summary/E03_SensorDepth.csv.gz}{\texttt{SensorDepth}} & 200 & 200 & soil, sensor depth & 4.94 & 199 & 86 & 583 & 86 & 146\\
E4 & \href{\repourl/blob/main/Data/Comprehensive_Summary/E04_Methods.csv.gz}{\texttt{Methods}} & 100 & 96 & soil, method & 5.01 & 164 & 88 & 542 & 83 & 147\\
E5 & \href{\repourl/blob/main/Data/Comprehensive_Summary/E05_SHP_Dynamics.csv.gz}{\texttt{SHP Dynamics}} & 300 & 300 & soil, method, tillage & 4.13 & 152 & 43 & 579 & 84 & 189\\
E6 & \href{\repourl/blob/main/Data/Comprehensive_Summary/E06_Mediterranean_Pollution.csv.gz}{\texttt{Mediterranean Pollution}} & 210 & 193 & soil, method, trigger, N rate & 3.64 & 84 & 239 & 620 & 62 & 118\\
E7 & \href{\repourl/blob/main/Data/Comprehensive_Summary/E07_Mulching.csv.gz}{\texttt{Mulching}} & 200 & 200 & soil, mulch & 5.29 & 183 & 96 & 555 & 74 & 161\\
E8 & \href{\repourl/blob/main/Data/Comprehensive_Summary/E08_Rain_Intensity_Moisture.csv.gz}{\texttt{Rain Intensity Moisture}} & 1000 & 1000 & initial suction, rain scale & 0.68 & 0 & -1 & 282 & -0 & 6\\
E9 & \href{\repourl/blob/main/Data/Comprehensive_Summary/E09_Soil_Moisture_Rain.csv.gz}{\texttt{Soil Moisture Rain}} & 1000 & 987 & soil, initial suction, rain scale & 0.75 & 0 & 19 & 247 & 33 & 7\\
E10 & \href{\repourl/blob/main/Data/Comprehensive_Summary/E10_Amount_Timing_Soil.csv.gz}{\texttt{Amount Timing Soil}} & 1000 & 954 & soil, rate, dose, window & 4.42 & 139 & 76 & 532 & 97 & 146\\
E11 & \href{\repourl/blob/main/Data/Comprehensive_Summary/E11_Temp_Evap_Rain.csv.gz}{\texttt{Temp Evap Rain}} & 1000 & 1000 & rain scale, warming, demand & 3.81 & 324 & 18 & 754 & 66 & 151\\
E12 & \href{\repourl/blob/main/Data/Comprehensive_Summary/E12_Sensor_Trigger_Rain.csv.gz}{\texttt{Sensor Trigger Rain}} & 1000 & 1000 & trigger, sensor depth, rain scale & 5.88 & 260 & 35 & 670 & 32 & 155\\
E13 & \href{\repourl/blob/main/Data/Comprehensive_Summary/E13_N_Water_Soil.csv.gz}{\texttt{N Water Soil}} & 1000 & 975 & soil, method, rate, dose, trigger, N rate & 4.00 & 162 & 83 & 547 & 79 & 207\\
E14 & \href{\repourl/blob/main/Data/Comprehensive_Summary/E14_Mulch_Rain.csv.gz}{\texttt{Mulch Rain}} & 1000 & 1000 & rain scale, mulch & 6.15 & 235 & 43 & 634 & 25 & 173\\
E15 & \href{\repourl/blob/main/Data/Comprehensive_Summary/E15_Method_Depth_Amount.csv.gz}{\texttt{Method Depth Amount}} & 1000 & 1000 & method, rate, dose, root depth & 6.02 & 217 & 61 & 619 & 37 & 157\\
E16 & \href{\repourl/blob/main/Data/Comprehensive_Summary/E16_Edge_Cases.csv.gz}{\texttt{Edge Cases}} & 770 & 740 & soil, method, rate, dose, trigger, sensor depth, rain scale, warming, N rate, demand & 2.17 & 344 & 185 & 600 & 93 & 224\\
\bottomrule
\end{tabular}
\end{table}

\begin{figure}[H]\centering
\includegraphics[width=\textwidth]{\figdir/data/fig_data_overview.pdf}
\caption{Yield, irrigation, drainage and nitrate leaching of the valid runs, by experiment.}\label{fig:data}
\end{figure}
\subsection*{out\_p4\_calibration.json}
\noindent\small Source: \nolinkurl{Manuscripts/reproducibility/p456/out_p4_calibration.json}; 140 values.\par
{\scriptsize\begin{longtable}{p{\dimexpr 0.4557\linewidth-2\tabcolsep\relax}p{\dimexpr 0.5443\linewidth-2\tabcolsep\relax}}\toprule\textbf{quantity} & \textbf{value}\\\midrule\endhead
{}truth.alphaF & 0.05303\\
{}truth.Kd & 0.02083\\
{}truth.alpha & 0.0001\\
{}loss0 & 1.409e-31\\
{}start & 0.204\\
{}runs[0].method & sequential\\
{}runs[0].seed & 1\\
{}runs[0].evals & 16\\
{}runs[0].loss & 0.0006235\\
{}runs[0].x.alphaF & 0.05\\
{}runs[0].x.Kd & 0.06\\
{}runs[0].x.alpha & 0.0003\\
{}runs[0].ms & 20\\
{}runs[0].err & 0.02553, 0.4594, 0.4771\\
{}runs[1].method & sequential\\
{}runs[1].seed & 2\\
{}runs[1].evals & 16\\
{}runs[1].loss & 0.0006235\\
{}runs[1].x.alphaF & 0.05\\
{}runs[1].x.Kd & 0.06\\
{}runs[1].x.alpha & 0.0003\\
{}runs[1].ms & 4\\
{}runs[1].err & 0.02553, 0.4594, 0.4771\\
{}runs[2].method & sequential\\
{}runs[2].seed & 3\\
{}runs[2].evals & 16\\
{}runs[2].loss & 0.0006235\\
{}runs[2].x.alphaF & 0.05\\
{}runs[2].x.Kd & 0.06\\
{}runs[2].x.alpha & 0.0003\\
{}runs[2].ms & 4\\
{}runs[2].err & 0.02553, 0.4594, 0.4771\\
{}runs[3].method & sequential\\
{}runs[3].seed & 4\\
{}runs[3].evals & 16\\
{}runs[3].loss & 0.0006235\\
{}runs[3].x.alphaF & 0.05\\
{}runs[3].x.Kd & 0.06\\
{}runs[3].x.alpha & 0.0003\\
{}runs[3].ms & 4\\
{}runs[3].err & 0.02553, 0.4594, 0.4771\\
{}runs[4].method & sequential\\
{}runs[4].seed & 5\\
{}runs[4].evals & 16\\
{}runs[4].loss & 0.0006235\\
{}runs[4].x.alphaF & 0.05\\
{}runs[4].x.Kd & 0.06\\
{}runs[4].x.alpha & 0.0003\\
{}runs[4].ms & 4\\
{}runs[4].err & 0.02553, 0.4594, 0.4771\\
{}runs[5].method & cma\\
{}runs[5].seed & 1\\
{}runs[5].evals & 154\\
{}runs[5].loss & 4.312e-06\\
{}runs[5].x.alphaF & 0.0533\\
{}runs[5].x.Kd & 0.004069\\
{}runs[5].x.alpha & 2.777e-05\\
{}runs[5].ms & 39\\
{}runs[5].err & 0.002251, 0.7093, 0.5564\\
{}runs[6].method & cma\\
{}runs[6].seed & 2\\
{}runs[6].evals & 154\\
{}runs[6].loss & 4.144e-06\\
{}runs[6].x.alphaF & 0.05313\\
{}runs[6].x.Kd & 0.006077\\
{}runs[6].x.alpha & 0.000431\\
{}runs[6].ms & 36\\
{}runs[6].err & 0.0008651, 0.535, 0.6345\\
{}runs[7].method & cma\\
{}runs[7].seed & 3\\
{}runs[7].evals & 154\\
{}runs[7].loss & 2.003e-06\\
{}runs[7].x.alphaF & 0.05286\\
{}runs[7].x.Kd & 0.05183\\
{}runs[7].x.alpha & 9.474e-05\\
{}runs[7].ms & 38\\
{}runs[7].err & 0.001402, 0.3958, 0.02349\\
{}runs[8].method & cma\\
{}runs[8].seed & 4\\
{}runs[8].evals & 154\\
{}runs[8].loss & 6.869e-06\\
{}runs[8].x.alphaF & 0.05337\\
{}runs[8].x.Kd & 0.02493\\
{}runs[8].x.alpha & 2.533e-05\\
{}runs[8].ms & 36\\
{}runs[8].err & 0.002804, 0.07801, 0.5964\\
{}runs[9].method & cma\\
{}runs[9].seed & 5\\
{}runs[9].evals & 154\\
{}runs[9].loss & 9.754e-09\\
{}runs[9].x.alphaF & 0.05301\\
{}runs[9].x.Kd & 0.003502\\
{}runs[9].x.alpha & 0.0001228\\
{}runs[9].ms & 37\\
{}runs[9].err & 0.0001234, 0.7744, 0.08912\\
{}runs[10].method & bayes\\
{}runs[10].seed & 1\\
{}runs[10].evals & 150\\
{}runs[10].loss & 6.888e-08\\
{}runs[10].x.alphaF & 0.05305\\
{}runs[10].x.Kd & 0.04165\\
{}runs[10].x.alpha & 2.105e-05\\
{}runs[10].ms & 1779\\
{}runs[10].err & 0.000177, 0.3008, 0.6768\\
{}runs[11].method & bayes\\
{}runs[11].seed & 2\\
{}runs[11].evals & 150\\
{}runs[11].loss & 7.706e-08\\
{}runs[11].x.alphaF & 0.05307\\
{}runs[11].x.Kd & 0.09905\\
{}runs[11].x.alpha & 6.412e-05\\
{}runs[11].ms & 1671\\
{}runs[11].err & 0.0003446, 0.6771, 0.193\\
{}runs[12].method & bayes\\
{}runs[12].seed & 3\\
{}runs[12].evals & 150\\
{}runs[12].loss & 8.24e-08\\
{}runs[12].x.alphaF & 0.05306\\
{}runs[12].x.Kd & 0.009166\\
{}runs[12].x.alpha & 1.795e-05\\
{}runs[12].ms & 1638\\
{}runs[12].err & 0.0003083, 0.3566, 0.746\\
{}runs[13].method & bayes\\
{}runs[13].seed & 4\\
{}runs[13].evals & 150\\
{}runs[13].loss & 1.321e-07\\
{}runs[13].x.alphaF & 0.05302\\
{}runs[13].x.Kd & 0.003837\\
{}runs[13].x.alpha & 2.652e-05\\
{}runs[13].ms & 1586\\
{}runs[13].err & 7.897e-05, 0.7348, 0.5764\\
{}runs[14].method & bayes\\
{}runs[14].seed & 5\\
{}runs[14].evals & 150\\
{}runs[14].loss & 2.571e-07\\
{}runs[14].x.alphaF & 0.053\\
{}runs[14].x.Kd & 0.05913\\
{}runs[14].x.alpha & 1.386e-05\\
{}runs[14].ms & 1666\\
{}runs[14].err & 0.000195, 0.453, 0.8582\\
\bottomrule\end{longtable}}

\subsection*{out\_p4\_design.json}
\noindent\small Source: \nolinkurl{Manuscripts/reproducibility/p456/out_p4_design.json}; 96 values.\par
{\scriptsize\begin{longtable}{p{\dimexpr 0.4557\linewidth-2\tabcolsep\relax}p{\dimexpr 0.5443\linewidth-2\tabcolsep\relax}}\toprule\textbf{quantity} & \textbf{value}\\\midrule\endhead
{}[0].id & E0\\
{}[0].name & Baselines\\
{}[0].gui & 9\\
{}[0].study & 10\\
{}[0].guiFactorCombos & 9\\
{}[0].ensCombos & 0\\
{}[1].id & E1\\
{}[1].name & Dose and duration\\
{}[1].gui & 17\\
{}[1].study & 200\\
{}[1].guiFactorCombos & 16\\
{}[1].ensCombos & 1\\
{}[2].id & E2\\
{}[2].name & Trigger pressure\\
{}[2].gui & 21\\
{}[2].study & 200\\
{}[2].guiFactorCombos & 18\\
{}[2].ensCombos & 200\\
{}[3].id & E3\\
{}[3].name & Sensor depth\\
{}[3].gui & 21\\
{}[3].study & 200\\
{}[3].guiFactorCombos & 18\\
{}[3].ensCombos & 200\\
{}[4].id & E4\\
{}[4].name & Methods\\
{}[4].gui & 6\\
{}[4].study & 100\\
{}[4].guiFactorCombos & 4\\
{}[4].ensCombos & 10\\
{}[5].id & E5\\
{}[5].name & Soil structure dynamics (SHP)\\
{}[5].gui & 72\\
{}[5].study & 300\\
{}[5].guiFactorCombos & 72\\
{}[5].ensCombos & 10\\
{}[6].id & E6\\
{}[6].name & Mediterranean pollution\\
{}[6].gui & 30\\
{}[6].study & 210\\
{}[6].guiFactorCombos & 30\\
{}[6].ensCombos & 10\\
{}[7].id & E7\\
{}[7].name & Mulching\\
{}[7].unsupported & True\\
{}[7].study & 200\\
{}[8].id & E8\\
{}[8].name & Rain intensity and moisture\\
{}[8].unsupported & True\\
{}[8].study & 1000\\
{}[9].id & E9\\
{}[9].name & Soil moisture and rain\\
{}[9].gui & 40\\
{}[9].study & 1000\\
{}[9].guiFactorCombos & 32\\
{}[9].ensCombos & 1000\\
{}[10].id & E10\\
{}[10].name & Amount, timing and soil\\
{}[10].gui & 20\\
{}[10].study & 1000\\
{}[10].guiFactorCombos & 18\\
{}[10].ensCombos & 100\\
{}[11].id & E11\\
{}[11].name & Temperature, evaporation, rain\\
{}[11].gui & 16\\
{}[11].study & 1000\\
{}[11].guiFactorCombos & 16\\
{}[11].ensCombos & 100\\
{}[12].id & E12\\
{}[12].name & Sensor, trigger, rain\\
{}[12].gui & 30\\
{}[12].study & 1000\\
{}[12].guiFactorCombos & 27\\
{}[12].ensCombos & 1000\\
{}[13].id & E13\\
{}[13].name & N, water, soil\\
{}[13].gui & 32\\
{}[13].study & 1000\\
{}[13].guiFactorCombos & 24\\
{}[13].ensCombos & 10\\
{}[14].id & E14\\
{}[14].name & Mulch and rain\\
{}[14].unsupported & True\\
{}[14].study & 1000\\
{}[15].id & E15\\
{}[15].name & Method, depth, amount\\
{}[15].gui & 30\\
{}[15].study & 1000\\
{}[15].guiFactorCombos & 27\\
{}[15].ensCombos & 100\\
{}[16].id & E16\\
{}[16].name & Edge cases\\
{}[16].gui & 24\\
{}[16].study & 770\\
{}[16].guiFactorCombos & 16\\
{}[16].ensCombos & 5\\
\bottomrule\end{longtable}}

\subsection*{out\_p4\_fidelity.csv}
\noindent\small Source: \nolinkurl{Manuscripts/reproducibility/p456/out_p4_fidelity.csv}; 51 rows. Workbook sheet \texttt{(CSV file)}.\par
{\scriptsize\setlength{\tabcolsep}{3pt}\begin{longtable}{>{\raggedright\arraybackslash}p{\dimexpr 0.1429\linewidth-2\tabcolsep\relax}>{\raggedright\arraybackslash}p{\dimexpr 0.1429\linewidth-2\tabcolsep\relax}>{\raggedright\arraybackslash}p{\dimexpr 0.1429\linewidth-2\tabcolsep\relax}>{\raggedright\arraybackslash}p{\dimexpr 0.1429\linewidth-2\tabcolsep\relax}>{\raggedright\arraybackslash}p{\dimexpr 0.1429\linewidth-2\tabcolsep\relax}>{\raggedright\arraybackslash}p{\dimexpr 0.1429\linewidth-2\tabcolsep\relax}>{\raggedright\arraybackslash}p{\dimexpr 0.1429\linewidth-2\tabcolsep\relax}}
\toprule
\textbf{exp} & \textbf{factor} & \textbf{resp} & \textbf{levels\_\allowbreak{}study} & \textbf{levels\_\allowbreak{}compact} & \textbf{rmse\_\allowbreak{}pct} & \textbf{max\_\allowbreak{}pct}\\\midrule\endhead
{}E2 & trig & yield\_Mg\_ha & 20 & 6 & 7.317 & 26.9\\
{}E2 & trig & irrigation\_mm & 20 & 6 & 6.443 & 22.11\\
{}E2 & trig & drainage\_1m\_mm & 20 & 6 & 8.957 & 36.94\\
{}E2 & trig & IE\_kg\_m3 & 20 & 6 & 8.133 & 28.55\\
{}E2 & trig & N\_leached\_kg\_ha & 20 & 6 & 7.183 & 23.64\\
{}E3 & z & yield\_Mg\_ha & 20 & 6 & 4.638 & 11.63\\
{}E3 & z & irrigation\_mm & 20 & 6 & 11.23 & 24.55\\
{}E3 & z & drainage\_1m\_mm & 20 & 6 & 4.563 & 9.679\\
{}E3 & z & IE\_kg\_m3 & 20 & 6 & 6.511 & 13.93\\
{}E3 & z & N\_leached\_kg\_ha & 20 & 6 & 4.849 & 13\\
{}E9 & trig & yield\_Mg\_ha & 10 & 4 & 11.68 & 33.78\\
{}E9 & trig & drainage\_1m\_mm & 10 & 4 & 26.19 & 79.45\\
{}E9 & trig & N\_leached\_kg\_ha & 10 & 4 & 16.2 & 49.6\\
{}E9 & P & yield\_Mg\_ha & 10 & 4 & 10.03 & 27.69\\
{}E9 & P & drainage\_1m\_mm & 10 & 4 & 12.32 & 35\\
{}E9 & P & N\_leached\_kg\_ha & 10 & 4 & 7.935 & 20.58\\
{}E10 & amt\_mm & yield\_Mg\_ha & 10 & 3 & 32.48 & 96.29\\
{}E10 & amt\_mm & irrigation\_mm & 10 & 3 & 29.85 & 88.58\\
{}E10 & amt\_mm & drainage\_1m\_mm & 10 & 3 & 36.08 & 88.19\\
{}E10 & amt\_mm & IE\_kg\_m3 & 10 & 3 & 27.91 & 85.16\\
{}E10 & amt\_mm & N\_leached\_kg\_ha & 10 & 3 & 30.38 & 90.76\\
{}E11 & dT & yield\_Mg\_ha & 10 & 4 & 8.061 & 19.75\\
{}E11 & dT & irrigation\_mm & 10 & 4 & 11.16 & 29.98\\
{}E11 & dT & drainage\_1m\_mm & 10 & 4 & 14.95 & 42.6\\
{}E11 & dT & IE\_kg\_m3 & 10 & 4 & 10.98 & 30.61\\
{}E11 & dT & N\_leached\_kg\_ha & 10 & 4 & 13.5 & 34.91\\
{}E11 & P & yield\_Mg\_ha & 10 & 4 & 6.598 & 14.63\\
{}E11 & P & irrigation\_mm & 10 & 4 & 9.053 & 24.73\\
{}E11 & P & drainage\_1m\_mm & 10 & 4 & 16.9 & 48.6\\
{}E11 & P & IE\_kg\_m3 & 10 & 4 & 10.36 & 28.78\\
{}E11 & P & N\_leached\_kg\_ha & 10 & 4 & 9.729 & 21.93\\
{}E12 & z & yield\_Mg\_ha & 10 & 3 & 11.16 & 23.65\\
{}E12 & z & irrigation\_mm & 10 & 3 & 13.24 & 24.2\\
{}E12 & z & drainage\_1m\_mm & 10 & 3 & 27.73 & 65.95\\
{}E12 & z & IE\_kg\_m3 & 10 & 3 & 21.02 & 50.32\\
{}E12 & z & N\_leached\_kg\_ha & 10 & 3 & 29.31 & 64.31\\
{}E12 & trig & yield\_Mg\_ha & 10 & 3 & 15.58 & 42.43\\
{}E12 & trig & irrigation\_mm & 10 & 3 & 13.2 & 28.82\\
{}E12 & trig & drainage\_1m\_mm & 10 & 3 & 17.55 & 52.43\\
{}E12 & trig & IE\_kg\_m3 & 10 & 3 & 18.63 & 54.2\\
{}E12 & trig & N\_leached\_kg\_ha & 10 & 3 & 16.13 & 44.94\\
{}E12 & P & yield\_Mg\_ha & 10 & 3 & 21.89 & 64.7\\
{}E12 & P & irrigation\_mm & 10 & 3 & 6.649 & 16.43\\
{}E12 & P & drainage\_1m\_mm & 10 & 3 & 8.067 & 23.93\\
{}E12 & P & IE\_kg\_m3 & 10 & 3 & 10.71 & 31.48\\
{}E12 & P & N\_leached\_kg\_ha & 10 & 3 & 30.16 & 89.24\\
{}E15 & root & yield\_Mg\_ha & 10 & 3 & 18.15 & 56.06\\
{}E15 & root & irrigation\_mm & 10 & 3 & 31.64 & 100\\
{}E15 & root & drainage\_1m\_mm & 10 & 3 & 27.96 & 88.29\\
{}E15 & root & IE\_kg\_m3 & 10 & 3 & 15.45 & 47.71\\
{}E15 & root & N\_leached\_kg\_ha & 10 & 3 & 14.68 & 44.85\\
\bottomrule
\end{longtable}}

\subsection*{out\_p4\_richards\_js.json}
\noindent\small Source: \nolinkurl{Manuscripts/reproducibility/p456/out_p4_richards_js.json}; 153 values.\par
{\scriptsize\begin{longtable}{p{\dimexpr 0.4557\linewidth-2\tabcolsep\relax}p{\dimexpr 0.5443\linewidth-2\tabcolsep\relax}}\toprule\textbf{quantity} & \textbf{value}\\\midrule\endhead
{}SOILS.sand.alpha & 0.145\\
{}SOILS.sand.n & 2.68\\
{}SOILS.sand.Ks & 29.7\\
{}SOILS.sand.Ths & 0.43\\
{}SOILS.sand.Thr & 0.045\\
{}SOILS.sand.l & 0.5\\
{}SOILS.loam.alpha & 0.036\\
{}SOILS.loam.n & 1.56\\
{}SOILS.loam.Ks & 1.04\\
{}SOILS.loam.Ths & 0.43\\
{}SOILS.loam.Thr & 0.078\\
{}SOILS.loam.l & 0.5\\
{}SOILS.siltclayloam.alpha & 0.01\\
{}SOILS.siltclayloam.n & 1.23\\
{}SOILS.siltclayloam.Ks & 0.07\\
{}SOILS.siltclayloam.Ths & 0.5\\
{}SOILS.siltclayloam.Thr & 0.089\\
{}SOILS.siltclayloam.l & 0.5\\
{}init.sand & 45.68\\
{}init.loam & 170.1\\
{}init.siltclayloam & 394.8\\
{}out.sand.1.infil & 200\\
{}out.sand.1.drain & 152.1\\
{}out.sand.1.runoff & 0\\
{}out.sand.1.store0 & None\\
{}out.sand.1.storeEnd & 93.56\\
{}out.sand.1.maxErr & 0\\
{}out.sand.1.sumErr & 0\\
{}out.sand.1.iters & 9964\\
{}out.sand.1.ms & 1586\\
{}out.sand.1.infilD & 200, 0, 0, 0, 0, 0\\
{}out.sand.1.drainD & 18.79, 96.26, 18.37, 8.962, 5.692, 4.046\\
{}out.sand.1.store & 226.9, 130.6, 112.3, 103.3, 97.6, 93.56\\
{}out.sand.2.infil & 200\\
{}out.sand.2.drain & 152.1\\
{}out.sand.2.runoff & 0\\
{}out.sand.2.store0 & None\\
{}out.sand.2.storeEnd & 93.56\\
{}out.sand.2.maxErr & 0\\
{}out.sand.2.sumErr & 0\\
{}out.sand.2.iters & 6306\\
{}out.sand.2.ms & 654\\
{}out.sand.2.infilD & 200, 0, 0, 0, 0, 0\\
{}out.sand.2.drainD & 18.79, 96.23, 18.37, 8.982, 5.7, 4.051\\
{}out.sand.2.store & 226.9, 130.7, 112.3, 103.3, 97.61, 93.56\\
{}out.sand.4.infil & 200\\
{}out.sand.4.drain & 152.2\\
{}out.sand.4.runoff & 0\\
{}out.sand.4.store0 & None\\
{}out.sand.4.storeEnd & 93.53\\
{}out.sand.4.maxErr & 0\\
{}out.sand.4.sumErr & 0\\
{}out.sand.4.iters & 4192\\
{}out.sand.4.ms & 298\\
{}out.sand.4.infilD & 200, 0, 0, 0, 0, 0\\
{}out.sand.4.drainD & 18.79, 96.16, 18.45, 9.001, 5.703, 4.054\\
{}out.sand.4.store & 226.9, 130.7, 112.3, 103.3, 97.58, 93.53\\
{}out.sand.0.5.infil & 200\\
{}out.sand.0.5.drain & 152.1\\
{}out.sand.0.5.runoff & 0\\
{}out.sand.0.5.store0 & None\\
{}out.sand.0.5.storeEnd & 93.54\\
{}out.sand.0.5.maxErr & 0\\
{}out.sand.0.5.sumErr & 0\\
{}out.sand.0.5.iters & 15784\\
{}out.sand.0.5.ms & 4958\\
{}out.sand.0.5.infilD & 200, 0, 0, 0, 0, 0\\
{}out.sand.0.5.drainD & 18.79, 96.27, 18.37, 8.99, 5.676, 4.047\\
{}out.sand.0.5.store & 226.9, 130.6, 112.2, 103.3, 97.59, 93.54\\
{}out.loam.1.infil & 200\\
{}out.loam.1.drain & 54.59\\
{}out.loam.1.runoff & 0\\
{}out.loam.1.store0 & None\\
{}out.loam.1.storeEnd & 315.5\\
{}out.loam.1.maxErr & 0\\
{}out.loam.1.sumErr & 0\\
{}out.loam.1.iters & 10486\\
{}out.loam.1.ms & 1599\\
{}out.loam.1.infilD & 200, 0, 0, 0, 0, 0\\
{}out.loam.1.drainD & 0.009, 5.772, 19.08, 12.88, 9.458, 7.396\\
{}out.loam.1.store & 370.1, 364.3, 345.2, 332.3, 322.9, 315.5\\
{}out.loam.2.infil & 200\\
{}out.loam.2.drain & 54.63\\
{}out.loam.2.runoff & 0\\
{}out.loam.2.store0 & None\\
{}out.loam.2.storeEnd & 315.4\\
{}out.loam.2.maxErr & 0\\
{}out.loam.2.sumErr & 0\\
{}out.loam.2.iters & 8549\\
{}out.loam.2.ms & 834\\
{}out.loam.2.infilD & 200, 0, 0, 0, 0, 0\\
{}out.loam.2.drainD & 0.009, 5.784, 19.04, 12.94, 9.454, 7.397\\
{}out.loam.2.store & 370.1, 364.3, 345.2, 332.3, 322.8, 315.4\\
{}out.loam.4.infil & 200\\
{}out.loam.4.drain & 54.58\\
{}out.loam.4.runoff & 0\\
{}out.loam.4.store0 & None\\
{}out.loam.4.storeEnd & 315.5\\
{}out.loam.4.maxErr & 0\\
{}out.loam.4.sumErr & 0\\
{}out.loam.4.iters & 7345\\
{}out.loam.4.ms & 477\\
{}out.loam.4.infilD & 200, 0, 0, 0, 0, 0\\
{}out.loam.4.drainD & 0.009, 5.825, 18.99, 12.88, 9.466, 7.402\\
{}out.loam.4.store & 370.1, 364.2, 345.2, 332.4, 322.9, 315.5\\
{}out.loam.0.5.infil & 200\\
{}out.loam.0.5.drain & 54.62\\
{}out.loam.0.5.runoff & 0\\
{}out.loam.0.5.store0 & None\\
{}out.loam.0.5.storeEnd & 315.4\\
{}out.loam.0.5.maxErr & 0\\
{}out.loam.0.5.sumErr & 0\\
{}out.loam.0.5.iters & 14781\\
{}out.loam.0.5.ms & 4257\\
{}out.loam.0.5.infilD & 200, 0, 0, 0, 0, 0\\
{}out.loam.0.5.drainD & 0.009, 5.773, 19.09, 12.9, 9.463, 7.389\\
{}out.loam.0.5.store & 370.1, 364.3, 345.2, 332.3, 322.8, 315.4\\
{}out.siltclayloam.1.infil & 21.77\\
{}out.siltclayloam.1.drain & 0.156\\
{}out.siltclayloam.1.runoff & 178.2\\
{}out.siltclayloam.1.store0 & None\\
{}out.siltclayloam.1.storeEnd & 416.4\\
{}out.siltclayloam.1.maxErr & 0.05\\
{}out.siltclayloam.1.sumErr & 0.05\\
{}out.siltclayloam.1.iters & 1345673\\
{}out.siltclayloam.1.ms & 182464\\
{}out.siltclayloam.1.infilD & 21.77, 0, 0, 0, 0, 0\\
{}out.siltclayloam.1.drainD & 0.026, 0.026, 0.026, 0.026, 0.026, 0.026\\
{}out.siltclayloam.1.store & 416.6, 416.6, 416.5, 416.5, 416.5, 416.4\\
{}out.siltclayloam.2.infil & 22.54\\
{}out.siltclayloam.2.drain & 0.156\\
{}out.siltclayloam.2.runoff & 177.5\\
{}out.siltclayloam.2.store0 & None\\
{}out.siltclayloam.2.storeEnd & 417.2\\
{}out.siltclayloam.2.maxErr & 0.074\\
{}out.siltclayloam.2.sumErr & 0.074\\
{}out.siltclayloam.2.iters & 266393\\
{}out.siltclayloam.2.ms & 20684\\
{}out.siltclayloam.2.infilD & 22.54, 0, 0, 0, 0, 0\\
{}out.siltclayloam.2.drainD & 0.026, 0.026, 0.026, 0.026, 0.026, 0.026\\
{}out.siltclayloam.2.store & 417.4, 417.3, 417.3, 417.3, 417.3, 417.2\\
{}out.siltclayloam.4.infil & 24.38\\
{}out.siltclayloam.4.drain & 0.156\\
{}out.siltclayloam.4.runoff & 175.6\\
{}out.siltclayloam.4.store0 & None\\
{}out.siltclayloam.4.storeEnd & 419\\
{}out.siltclayloam.4.maxErr & 0.007\\
{}out.siltclayloam.4.sumErr & -0.007\\
{}out.siltclayloam.4.iters & 385896\\
{}out.siltclayloam.4.ms & 21266\\
{}out.siltclayloam.4.infilD & 24.38, 0, 0, 0, 0, 0\\
{}out.siltclayloam.4.drainD & 0.026, 0.026, 0.026, 0.026, 0.026, 0.026\\
{}out.siltclayloam.4.store & 419.1, 419.1, 419.1, 419, 419, 419\\
\bottomrule\end{longtable}}

\subsection*{out\_p4\_richards\_ref.json}
\noindent\small Source: \nolinkurl{Manuscripts/reproducibility/p456/out_p4_richards_ref.json}; 8 values.\par
{\scriptsize\begin{longtable}{p{\dimexpr 0.4557\linewidth-2\tabcolsep\relax}p{\dimexpr 0.5443\linewidth-2\tabcolsep\relax}}\toprule\textbf{quantity} & \textbf{value}\\\midrule\endhead
{}sand.infil & 200\\
{}sand.drain & 152.6\\
{}sand.storeEnd & 93.12\\
{}sand.S0 & 45.68\\
{}loam.infil & 200\\
{}loam.drain & 55.12\\
{}loam.storeEnd & 314.9\\
{}loam.S0 & 170.1\\
\bottomrule\end{longtable}}

\subsection*{out\_p4\_runtime.csv}
\noindent\small Source: \nolinkurl{Manuscripts/reproducibility/p456/out_p4_runtime.csv}; 17 rows. Workbook sheet \texttt{(CSV file)}.\par
{\scriptsize\setlength{\tabcolsep}{3pt}\begin{longtable}{>{\raggedright\arraybackslash}p{\dimexpr 0.1667\linewidth-2\tabcolsep\relax}>{\raggedright\arraybackslash}p{\dimexpr 0.1667\linewidth-2\tabcolsep\relax}>{\raggedright\arraybackslash}p{\dimexpr 0.1667\linewidth-2\tabcolsep\relax}>{\raggedright\arraybackslash}p{\dimexpr 0.1667\linewidth-2\tabcolsep\relax}>{\raggedright\arraybackslash}p{\dimexpr 0.1667\linewidth-2\tabcolsep\relax}>{\raggedright\arraybackslash}p{\dimexpr 0.1667\linewidth-2\tabcolsep\relax}}
\toprule
\textbf{exp} & \textbf{name} & \textbf{study} & \textbf{study\_\allowbreak{}hours} & \textbf{gui} & \textbf{gui\_\allowbreak{}hours}\\\midrule\endhead
{}E0 & Baselines & 10 & 0.6411 & 9 & 0.5925\\
{}E1 & Dose and duration & 200 & 10.99 & 17 & 0.8878\\
{}E2 & Trigger pressure & 200 & 7.731 & 21 & 0.805\\
{}E3 & Sensor depth & 200 & 8.099 & 21 & 0.84\\
{}E4 & Methods & 100 & 4.1 & 6 & 0.2408\\
{}E5 & Soil structure dynamics (SHP) & 300 & 15.77 & 72 & 3.45\\
{}E6 & Mediterranean pollution & 210 & 6.937 & 30 & 0.7833\\
{}E7 & Mulching & 200 & 8.949 & nan & nan\\
{}E8 & Rain intensity and moisture & 1000 & 1.727 & nan & nan\\
{}E9 & Soil moisture and rain & 1000 & 1.957 & 40 & 0.07778\\
{}E10 & Amount, timing and soil & 1000 & 70.39 & 20 & 0.7111\\
{}E11 & Temperature, evaporation, rain & 1000 & 41.98 & 16 & 0.6578\\
{}E12 & Sensor, trigger, rain & 1000 & 42.92 & 30 & 1.225\\
{}E13 & N, water, soil & 1000 & 57.44 & 32 & 1.929\\
{}E14 & Mulch and rain & 1000 & 48.01 & nan & nan\\
{}E15 & Method, depth, amount & 1000 & 43.58 & 30 & 1.275\\
{}E16 & Edge cases & 770 & 48.22 & 24 & 1.21\\
\bottomrule
\end{longtable}}

\subsection*{out\_p5\_bridge.json}
\noindent\small Source: \nolinkurl{Manuscripts/reproducibility/p456/out_p5_bridge.json}; 20 values.\par
{\scriptsize\begin{longtable}{p{\dimexpr 0.4557\linewidth-2\tabcolsep\relax}p{\dimexpr 0.5443\linewidth-2\tabcolsep\relax}}\toprule\textbf{quantity} & \textbf{value}\\\midrule\endhead
{}parser[0] & annual-field-nitrogen-3b.dlf, 11, 11, 2.246e-16\\
{}parser[1] & annual-field-nitrogen-4b.dlf, 11, 11, 1.703e-16\\
{}parser[2] & annual-tracer-3b.dlf, 11, 11, 7.006e-34\\
{}parser[3] & daily-water-flux.dlf, 26, 26, 3.277e-16\\
{}parser[4] & optim-field-nitrogen.dlf, 1, 1, 0\\
{}parser[5] & soil-chemical-110cm.dlf, 11665, 11665, 2.918e-15\\
{}n.annual-field-nitrogen-3b.dlf.Nfert & 0\\
{}n.annual-field-nitrogen-3b.dlf.Nup & 0.045\\
{}n.annual-field-nitrogen-3b.dlf.NL & 0.00478\\
{}n.annual-field-nitrogen-3b.dlf.Nmin & 0.049\\
{}n.annual-field-nitrogen-3b.dlf.Nden & 0.00473\\
{}n.annual-field-nitrogen-3b.dlf.maxabs\_Error\_kgN & 0\\
{}n.annual-field-nitrogen-3b.dlf.seasons & 10\\
{}n.annual-field-nitrogen-4b.dlf.Nfert & 0\\
{}n.annual-field-nitrogen-4b.dlf.Nup & 0.049\\
{}n.annual-field-nitrogen-4b.dlf.NL & 0.0041\\
{}n.annual-field-nitrogen-4b.dlf.Nmin & 0.048\\
{}n.annual-field-nitrogen-4b.dlf.Nden & 0.00442\\
{}n.annual-field-nitrogen-4b.dlf.maxabs\_Error\_kgN & 0\\
{}n.annual-field-nitrogen-4b.dlf.seasons & 10\\
\bottomrule\end{longtable}}

\subsection*{out\_p5\_bridge\_js.json}
\noindent\small Source: \nolinkurl{Manuscripts/reproducibility/p456/out_p5_bridge_js.json}; 832 values, first 220 shown.\par
{\scriptsize\begin{longtable}{p{\dimexpr 0.4557\linewidth-2\tabcolsep\relax}p{\dimexpr 0.5443\linewidth-2\tabcolsep\relax}}\toprule\textbf{quantity} & \textbf{value}\\\midrule\endhead
{}deniArr[0] & 3.206e-11, 3.206e-11, 1.88e-11, 4.709e-11, 4.709e-11, 2.869e-11\\
{}deniArr[1] & 3.83e-11, 3.83e-11, 2.235e-11, 5.479e-11, 5.479e-11, 3.35e-11\\
{}deniArr[2] & 4.481e-11, 4.481e-11, 2.599e-11, 6.21e-11, 6.21e-11, 3.816e-11\\
{}deniArr[3] & 3.869e-11, 3.869e-11, 2.239e-11, 5.305e-11, 5.305e-11, 3.265e-11\\
{}deniArr[4] & 2.408e-11, 2.408e-11, 1.396e-11, 3.312e-11, 3.312e-11, 2.037e-11\\
{}deniArr[5] & 5.238e-12, 5.238e-12, 3.111e-12, 1.273e-11, 1.273e-11, 7.671e-12\\
{}deniArr[6] & 2.994e-12, 2.994e-12, 1.785e-12, 7.394e-12, 7.394e-12, 4.447e-12\\
{}deniArr[7] & 1.63e-12, 1.63e-12, 9.747e-13, 4.065e-12, 4.065e-12, 2.442e-12\\
{}deniArr[8] & 8.865e-13, 8.865e-13, 5.309e-13, 2.222e-12, 2.222e-12, 1.334e-12\\
{}deniArr[9] & 4.67e-13, 4.67e-13, 2.799e-13, 1.174e-12, 1.174e-12, 7.048e-13\\
{}deniArr[10] & 2.463e-13, 2.463e-13, 1.477e-13, 6.204e-13, 6.204e-13, 3.723e-13\\
{}deniArr[11] & 1.221e-13, 1.221e-13, 7.327e-14, 3.079e-13, 3.079e-13, 1.848e-13\\
{}deniArr[12] & 2.007e-13, 2.007e-13, 1.204e-13, 2.665e-13, 2.665e-13, 1.599e-13\\
{}deniArr[13] & 9.756e-14, 9.756e-14, 5.853e-14, 1.295e-13, 1.295e-13, 7.772e-14\\
{}deniArr[14] & 0, 0, 0, 0, 0, 0\\
{}deniArr[15] & 0, 0, 0, 0, 0, 0\\
{}deniArr[16] & 0, 0, 0, 0, 0, 0\\
{}deniArr[17] & 0, 0, 0, 0, 0, 0\\
{}deniArr[18] & 0, 0, 0, 0, 0, 0\\
{}deniArr[19] & 0, 0, 0, 0, 0, 0\\
{}deniArr[20] & 0, 0, 0, 0, 0, 0\\
{}deniArr[21] & 0, 0, 0, 0, 0, 0\\
{}deniArr[22] & 0, 0, 0, 0, 0, 0\\
{}deniArr[23] & 0, 0, 0, 0, 0, 0\\
{}deniArr[24] & 0, 0, 0, 0, 0, 0\\
{}deniArr[25] & 0, 0, 0, 0, 0, 0\\
{}deniArr[26] & 0, 0, 0, 0, 0, 0\\
{}deniArr[27] & 0, 0, 0, 0, 0, 0\\
{}deniArr[28] & 0, 0, 0, 0, 0, 0\\
{}deniArr[29] & 0, 0, 0, 0, 0, 0\\
{}deniArr[30] & 0, 0, 0, 0, 0, 0\\
{}deniArr[31] & 0, 0, 0, 0, 0, 0\\
{}deniArr[32] & 0, 0, 0, 0, 0, 0\\
{}deniArr[33] & 1.124e-10, 1.124e-10, 6.329e-11, 1.363e-10, 1.363e-10, 8.598e-11\\
{}deniArr[34] & 7.183e-11, 7.183e-11, 4.113e-11, 9.431e-11, 9.431e-11, 5.856e-11\\
{}deniArr[35] & 5.099e-11, 5.099e-11, 2.946e-11, 6.949e-11, 6.949e-11, 4.283e-11\\
{}deniArr[36] & 3.823e-11, 3.823e-11, 2.213e-11, 5.249e-11, 5.249e-11, 3.23e-11\\
{}deniArr[37] & 2.373e-11, 2.373e-11, 1.376e-11, 3.264e-11, 3.264e-11, 2.006e-11\\
{}deniArr[38] & 5.254e-12, 5.254e-12, 3.12e-12, 1.276e-11, 1.276e-11, 7.687e-12\\
{}deniArr[39] & 2.973e-12, 2.973e-12, 1.773e-12, 7.343e-12, 7.343e-12, 4.417e-12\\
{}deniArr[40] & 1.634e-12, 1.634e-12, 9.773e-13, 4.075e-12, 4.075e-12, 2.448e-12\\
{}deniArr[41] & 8.821e-13, 8.821e-13, 5.283e-13, 2.211e-12, 2.211e-12, 1.328e-12\\
{}deniArr[42] & 4.679e-13, 4.679e-13, 2.805e-13, 1.177e-12, 1.177e-12, 7.062e-13\\
{}deniArr[43] & 2.457e-13, 2.457e-13, 1.474e-13, 6.189e-13, 6.189e-13, 3.714e-13\\
{}deniArr[44] & 1.221e-13, 1.221e-13, 7.323e-14, 3.079e-13, 3.079e-13, 1.848e-13\\
{}deniArr[45] & 2.007e-13, 2.007e-13, 1.204e-13, 2.666e-13, 2.666e-13, 1.6e-13\\
{}deniArr[46] & 9.747e-14, 9.747e-14, 5.848e-14, 1.294e-13, 1.294e-13, 7.765e-14\\
{}deniArr[47] & 0, 0, 0, 0, 0, 0\\
{}deniArr[48] & 0, 0, 0, 0, 0, 0\\
{}deniArr[49] & 0, 0, 0, 0, 0, 0\\
{}deniArr[50] & 0, 0, 0, 0, 0, 0\\
{}deniArr[51] & 0, 0, 0, 0, 0, 0\\
{}deniArr[52] & 0, 0, 0, 0, 0, 0\\
{}deniArr[53] & 0, 0, 0, 0, 0, 0\\
{}deniArr[54] & 0, 0, 0, 0, 0, 0\\
{}deniArr[55] & 0, 0, 0, 0, 0, 0\\
{}deniArr[56] & 0, 0, 0, 0, 0, 0\\
{}deniArr[57] & 0, 0, 0, 0, 0, 0\\
{}deniArr[58] & 0, 0, 0, 0, 0, 0\\
{}deniArr[59] & 0, 0, 0, 0, 0, 0\\
{}deniArr[60] & 0, 0, 0, 0, 0, 0\\
{}deniArr[61] & 0, 0, 0, 0, 0, 0\\
{}deniArr[62] & 0, 0, 0, 0, 0, 0\\
{}deniArr[63] & 0, 0, 0, 0, 0, 0\\
{}deniArr[64] & 0, 0, 0, 0, 0, 0\\
{}deniArr[65] & 0, 0, 0, 0, 0, 0\\
{}deniArr[66] & 1.964e-10, 1.964e-10, 1.074e-10, 2.027e-10, 2.027e-10, 1.32e-10\\
{}deniArr[67] & 1.222e-10, 1.222e-10, 6.842e-11, 1.432e-10, 1.432e-10, 9.084e-11\\
{}deniArr[68] & 6.988e-11, 6.988e-11, 3.993e-11, 9.051e-11, 9.051e-11, 5.63e-11\\
{}deniArr[69] & 4.024e-11, 4.024e-11, 2.326e-11, 5.486e-11, 5.486e-11, 3.379e-11\\
{}deniArr[70] & 2.371e-11, 2.371e-11, 1.376e-11, 3.255e-11, 3.255e-11, 2e-11\\
{}deniArr[71] & 5.244e-12, 5.244e-12, 3.114e-12, 1.272e-11, 1.272e-11, 7.665e-12\\
{}deniArr[72] & 2.963e-12, 2.963e-12, 1.767e-12, 7.321e-12, 7.321e-12, 4.403e-12\\
{}deniArr[73] & 1.632e-12, 1.632e-12, 9.761e-13, 4.07e-12, 4.07e-12, 2.445e-12\\
{}deniArr[74] & 8.806e-13, 8.806e-13, 5.274e-13, 2.208e-12, 2.208e-12, 1.325e-12\\
{}deniArr[75] & 4.677e-13, 4.677e-13, 2.803e-13, 1.176e-12, 1.176e-12, 7.058e-13\\
{}deniArr[76] & 2.455e-13, 2.455e-13, 1.472e-13, 6.183e-13, 6.183e-13, 3.711e-13\\
{}deniArr[77] & 1.22e-13, 1.22e-13, 7.318e-14, 3.077e-13, 3.077e-13, 1.847e-13\\
{}deniArr[78] & 2.006e-13, 2.006e-13, 1.204e-13, 2.666e-13, 2.666e-13, 1.6e-13\\
{}deniArr[79] & 9.741e-14, 9.741e-14, 5.844e-14, 1.293e-13, 1.293e-13, 7.76e-14\\
{}deniArr[80] & 0, 0, 0, 0, 0, 0\\
{}deniArr[81] & 0, 0, 0, 0, 0, 0\\
{}deniArr[82] & 0, 0, 0, 0, 0, 0\\
{}deniArr[83] & 0, 0, 0, 0, 0, 0\\
{}deniArr[84] & 0, 0, 0, 0, 0, 0\\
{}deniArr[85] & 0, 0, 0, 0, 0, 0\\
{}deniArr[86] & 0, 0, 0, 0, 0, 0\\
{}deniArr[87] & 0, 0, 0, 0, 0, 0\\
{}deniArr[88] & 0, 0, 0, 0, 0, 0\\
{}deniArr[89] & 0, 0, 0, 0, 0, 0\\
{}deniArr[90] & 0, 0, 0, 0, 0, 0\\
{}deniArr[91] & 0, 0, 0, 0, 0, 0\\
{}deniArr[92] & 0, 0, 0, 0, 0, 0\\
{}deniArr[93] & 0, 0, 0, 0, 0, 0\\
{}deniArr[94] & 0, 0, 0, 0, 0, 0\\
{}deniArr[95] & 0, 0, 0, 0, 0, 0\\
{}deniArr[96] & 0, 0, 0, 0, 0, 0\\
{}deniArr[97] & 0, 0, 0, 0, 0, 0\\
{}deniArr[98] & 0, 0, 0, 0, 0, 0\\
{}deniArr[99] & 2.86e-10, 2.86e-10, 1.53e-10, 2.548e-10, 2.548e-10, 1.715e-10\\
{}deniArr[100] & 1.737e-10, 1.737e-10, 9.543e-11, 1.836e-10, 1.836e-10, 1.189e-10\\
{}deniArr[101] & 9.365e-11, 9.365e-11, 5.287e-11, 1.142e-10, 1.142e-10, 7.186e-11\\
{}deniArr[102] & 4.498e-11, 4.498e-11, 2.591e-11, 6.029e-11, 6.029e-11, 3.725e-11\\
{}deniArr[103] & 2.415e-11, 2.415e-11, 1.4e-11, 3.3e-11, 3.3e-11, 2.029e-11\\
{}deniArr[104] & 5.245e-12, 5.245e-12, 3.115e-12, 1.271e-11, 1.271e-11, 7.659e-12\\
{}deniArr[105] & 2.956e-12, 2.956e-12, 1.763e-12, 7.303e-12, 7.303e-12, 4.392e-12\\
{}deniArr[106] & 1.629e-12, 1.629e-12, 9.744e-13, 4.063e-12, 4.063e-12, 2.441e-12\\
{}deniArr[107] & 8.794e-13, 8.794e-13, 5.267e-13, 2.205e-12, 2.205e-12, 1.324e-12\\
{}deniArr[108] & 4.673e-13, 4.673e-13, 2.801e-13, 1.175e-12, 1.175e-12, 7.052e-13\\
{}deniArr[109] & 2.454e-13, 2.454e-13, 1.471e-13, 6.179e-13, 6.179e-13, 3.708e-13\\
{}deniArr[110] & 1.219e-13, 1.219e-13, 7.312e-14, 3.076e-13, 3.076e-13, 1.846e-13\\
{}deniArr[111] & 2.004e-13, 2.004e-13, 1.203e-13, 2.664e-13, 2.664e-13, 1.599e-13\\
{}deniArr[112] & 9.735e-14, 9.735e-14, 5.841e-14, 1.293e-13, 1.293e-13, 7.756e-14\\
{}deniArr[113] & 0, 0, 0, 0, 0, 0\\
{}deniArr[114] & 0, 0, 0, 0, 0, 0\\
{}deniArr[115] & 0, 0, 0, 0, 0, 0\\
{}deniArr[116] & 0, 0, 0, 0, 0, 0\\
{}deniArr[117] & 0, 0, 0, 0, 0, 0\\
{}deniArr[118] & 0, 0, 0, 0, 0, 0\\
{}deniArr[119] & 0, 0, 0, 0, 0, 0\\
{}deniArr[120] & 0, 0, 0, 0, 0, 0\\
{}deniArr[121] & 0, 0, 0, 0, 0, 0\\
{}deniArr[122] & 0, 0, 0, 0, 0, 0\\
{}deniArr[123] & 0, 0, 0, 0, 0, 0\\
{}deniArr[124] & 0, 0, 0, 0, 0, 0\\
{}deniArr[125] & 0, 0, 0, 0, 0, 0\\
{}deniArr[126] & 0, 0, 0, 0, 0, 0\\
{}deniArr[127] & 0, 0, 0, 0, 0, 0\\
{}deniArr[128] & 0, 0, 0, 0, 0, 0\\
{}deniArr[129] & 0, 0, 0, 0, 0, 0\\
{}deniArr[130] & 0, 0, 0, 0, 0, 0\\
{}deniArr[131] & 0, 0, 0, 0, 0, 0\\
{}deniArr[132] & 3.764e-10, 3.764e-10, 1.979e-10, 2.94e-10, 2.94e-10, 2.043e-10\\
{}deniArr[133] & 2.296e-10, 2.296e-10, 1.242e-10, 2.196e-10, 2.196e-10, 1.453e-10\\
{}deniArr[134] & 1.196e-10, 1.196e-10, 6.675e-11, 1.373e-10, 1.373e-10, 8.741e-11\\
{}deniArr[135] & 5.156e-11, 5.156e-11, 2.956e-11, 6.756e-11, 6.756e-11, 4.191e-11\\
{}deniArr[136] & 2.508e-11, 2.508e-11, 1.453e-11, 3.405e-11, 3.405e-11, 2.095e-11\\
{}deniArr[137] & 5.274e-12, 5.274e-12, 3.132e-12, 1.277e-11, 1.277e-11, 7.693e-12\\
{}deniArr[138] & 2.952e-12, 2.952e-12, 1.761e-12, 7.293e-12, 7.293e-12, 4.386e-12\\
{}deniArr[139] & 1.627e-12, 1.627e-12, 9.728e-13, 4.057e-12, 4.057e-12, 2.437e-12\\
{}deniArr[140] & 8.783e-13, 8.783e-13, 5.26e-13, 2.202e-12, 2.202e-12, 1.322e-12\\
{}deniArr[141] & 4.669e-13, 4.669e-13, 2.799e-13, 1.174e-12, 1.174e-12, 7.046e-13\\
{}deniArr[142] & 2.452e-13, 2.452e-13, 1.47e-13, 6.175e-13, 6.175e-13, 3.706e-13\\
{}deniArr[143] & 1.218e-13, 1.218e-13, 7.306e-14, 3.074e-13, 3.074e-13, 1.845e-13\\
{}deniArr[144] & 2.003e-13, 2.003e-13, 1.201e-13, 2.663e-13, 2.663e-13, 1.598e-13\\
{}deniArr[145] & 9.729e-14, 9.729e-14, 5.837e-14, 1.292e-13, 1.292e-13, 7.751e-14\\
{}deniArr[146] & 0, 0, 0, 0, 0, 0\\
{}deniArr[147] & 0, 0, 0, 0, 0, 0\\
{}deniArr[148] & 0, 0, 0, 0, 0, 0\\
{}deniArr[149] & 0, 0, 0, 0, 0, 0\\
{}deniArr[150] & 0, 0, 0, 0, 0, 0\\
{}deniArr[151] & 0, 0, 0, 0, 0, 0\\
{}deniArr[152] & 0, 0, 0, 0, 0, 0\\
{}deniArr[153] & 0, 0, 0, 0, 0, 0\\
{}deniArr[154] & 0, 0, 0, 0, 0, 0\\
{}deniArr[155] & 0, 0, 0, 0, 0, 0\\
{}deniArr[156] & 0, 0, 0, 0, 0, 0\\
{}deniArr[157] & 0, 0, 0, 0, 0, 0\\
{}deniArr[158] & 0, 0, 0, 0, 0, 0\\
{}deniArr[159] & 0, 0, 0, 0, 0, 0\\
{}deniArr[160] & 0, 0, 0, 0, 0, 0\\
{}deniArr[161] & 0, 0, 0, 0, 0, 0\\
{}deniArr[162] & 0, 0, 0, 0, 0, 0\\
{}deniArr[163] & 0, 0, 0, 0, 0, 0\\
{}deniArr[164] & 0, 0, 0, 0, 0, 0\\
{}deniArr[165] & 4.711e-10, 4.711e-10, 2.444e-10, 3.245e-10, 3.245e-10, 2.33e-10\\
{}deniArr[166] & 2.868e-10, 2.868e-10, 1.53e-10, 2.497e-10, 2.497e-10, 1.689e-10\\
{}deniArr[167] & 1.476e-10, 1.476e-10, 8.148e-11, 1.594e-10, 1.594e-10, 1.027e-10\\
{}deniArr[168] & 5.941e-11, 5.941e-11, 3.389e-11, 7.585e-11, 7.585e-11, 4.727e-11\\
{}deniArr[169] & 2.648e-11, 2.648e-11, 1.532e-11, 3.562e-11, 3.562e-11, 2.194e-11\\
{}deniArr[170] & 5.342e-12, 5.342e-12, 3.172e-12, 1.291e-11, 1.291e-11, 7.781e-12\\
{}deniArr[171] & 2.954e-12, 2.954e-12, 1.762e-12, 7.297e-12, 7.297e-12, 4.389e-12\\
{}deniArr[172] & 1.625e-12, 1.625e-12, 9.717e-13, 4.052e-12, 4.052e-12, 2.434e-12\\
{}deniArr[173] & 8.773e-13, 8.773e-13, 5.254e-13, 2.199e-12, 2.199e-12, 1.321e-12\\
{}deniArr[174] & 4.665e-13, 4.665e-13, 2.796e-13, 1.173e-12, 1.173e-12, 7.04e-13\\
{}deniArr[175] & 2.45e-13, 2.45e-13, 1.47e-13, 6.171e-13, 6.171e-13, 3.703e-13\\
{}deniArr[176] & 1.217e-13, 1.217e-13, 7.299e-14, 3.072e-13, 3.072e-13, 1.843e-13\\
{}deniArr[177] & 2.001e-13, 2.001e-13, 1.2e-13, 2.661e-13, 2.661e-13, 1.597e-13\\
{}deniArr[178] & 9.722e-14, 9.722e-14, 5.833e-14, 1.291e-13, 1.291e-13, 7.746e-14\\
{}deniArr[179] & 0, 0, 0, 0, 0, 0\\
{}deniArr[180] & 0, 0, 0, 0, 0, 0\\
{}deniArr[181] & 0, 0, 0, 0, 0, 0\\
{}deniArr[182] & 0, 0, 0, 0, 0, 0\\
{}deniArr[183] & 0, 0, 0, 0, 0, 0\\
{}deniArr[184] & 0, 0, 0, 0, 0, 0\\
{}deniArr[185] & 0, 0, 0, 0, 0, 0\\
{}deniArr[186] & 0, 0, 0, 0, 0, 0\\
{}deniArr[187] & 0, 0, 0, 0, 0, 0\\
{}deniArr[188] & 0, 0, 0, 0, 0, 0\\
{}deniArr[189] & 0, 0, 0, 0, 0, 0\\
{}deniArr[190] & 0, 0, 0, 0, 0, 0\\
{}deniArr[191] & 0, 0, 0, 0, 0, 0\\
{}deniArr[192] & 0, 0, 0, 0, 0, 0\\
{}deniArr[193] & 0, 0, 0, 0, 0, 0\\
{}deniArr[194] & 0, 0, 0, 0, 0, 0\\
{}deniArr[195] & 0, 0, 0, 0, 0, 0\\
{}deniArr[196] & 0, 0, 0, 0, 0, 0\\
{}deniArr[197] & 0, 0, 0, 0, 0, 0\\
{}deniArr[198] & 5.655e-10, 5.655e-10, 2.901e-10, 3.469e-10, 3.469e-10, 2.573e-10\\
{}deniArr[199] & 3.463e-10, 3.463e-10, 1.826e-10, 2.755e-10, 2.755e-10, 1.904e-10\\
{}deniArr[200] & 1.769e-10, 1.769e-10, 9.67e-11, 1.8e-10, 1.8e-10, 1.175e-10\\
{}deniArr[201] & 6.822e-11, 6.822e-11, 3.87e-11, 8.468e-11, 8.468e-11, 5.304e-11\\
{}deniArr[202] & 2.827e-11, 2.827e-11, 1.632e-11, 3.762e-11, 3.762e-11, 2.321e-11\\
{}deniArr[203] & 5.45e-12, 5.45e-12, 3.236e-12, 1.315e-11, 1.315e-11, 7.925e-12\\
{}deniArr[204] & 2.964e-12, 2.964e-12, 1.768e-12, 7.322e-12, 7.322e-12, 4.404e-12\\
{}deniArr[205] & 1.624e-12, 1.624e-12, 9.713e-13, 4.05e-12, 4.05e-12, 2.433e-12\\
{}deniArr[206] & 8.764e-13, 8.764e-13, 5.249e-13, 2.197e-12, 2.197e-12, 1.319e-12\\
{}deniArr[207] & 4.661e-13, 4.661e-13, 2.794e-13, 1.172e-12, 1.172e-12, 7.034e-13\\
{}deniArr[208] & 2.449e-13, 2.449e-13, 1.469e-13, 6.167e-13, 6.167e-13, 3.701e-13\\
{}deniArr[209] & 1.216e-13, 1.216e-13, 7.292e-14, 3.07e-13, 3.07e-13, 1.842e-13\\
{}deniArr[210] & 1.998e-13, 1.998e-13, 1.199e-13, 2.659e-13, 2.659e-13, 1.596e-13\\
{}deniArr[211] & 9.715e-14, 9.715e-14, 5.829e-14, 1.29e-13, 1.29e-13, 7.741e-14\\
{}deniArr[212] & 0, 0, 0, 0, 0, 0\\
{}deniArr[213] & 0, 0, 0, 0, 0, 0\\
{}deniArr[214] & 0, 0, 0, 0, 0, 0\\
{}deniArr[215] & 0, 0, 0, 0, 0, 0\\
{}deniArr[216] & 0, 0, 0, 0, 0, 0\\
{}deniArr[217] & 0, 0, 0, 0, 0, 0\\
{}deniArr[218] & 0, 0, 0, 0, 0, 0\\
{}deniArr[219] & 0, 0, 0, 0, 0, 0\\
\bottomrule\end{longtable}}

\subsection*{out\_p5\_chem.json}
\noindent\small Source: \nolinkurl{Manuscripts/reproducibility/p456/out_p5_chem.json}; 6 values.\par
{\scriptsize\begin{longtable}{p{\dimexpr 0.4557\linewidth-2\tabcolsep\relax}p{\dimexpr 0.5443\linewidth-2\tabcolsep\relax}}\toprule\textbf{quantity} & \textbf{value}\\\midrule\endhead
{}resid & -0.1097, 0.1082, 0.0011, -3.553e-15, 0.0007, -0.1228, 0.1221, 0, -0.0002, 0.001\\
{}maxAbs & 0.1228\\
{}sumAbs & 0.4658\\
{}applied & 2290\\
{}leachTotal & 393\\
{}uptakeTotal & 1934\\
\bottomrule\end{longtable}}

\subsection*{out\_p6\_crops.json}
\noindent\small Source: \nolinkurl{Manuscripts/reproducibility/p456/out_p6_crops.json}; 1495 values, first 220 shown.\par
{\scriptsize\begin{longtable}{p{\dimexpr 0.4557\linewidth-2\tabcolsep\relax}p{\dimexpr 0.5443\linewidth-2\tabcolsep\relax}}\toprule\textbf{quantity} & \textbf{value}\\\midrule\endhead
{}n & 48\\
{}groups.field & 19\\
{}groups.perennial & 29\\
{}habits.annual & 19\\
{}habits.deciduous & 18\\
{}habits.evergreen & 11\\
{}combos & 110\\
{}issues[0][0] & citrus\\
{}issues[0][1] & Young\\
{}issues[0][2] & Kc mid < Kc ini\\
{}issues[1][0] & citrus\\
{}issues[1][1] & Low density\\
{}issues[1][2] & Kc mid < Kc ini\\
{}issues[2][0] & citrus\\
{}issues[2][1] & Medium density, short\\
{}issues[2][2] & Kc mid < Kc ini\\
{}issues[3][0] & citrus\\
{}issues[3][1] & Medium density, tall\\
{}issues[3][2] & Kc mid < Kc ini\\
{}issues[4][0] & citrus\\
{}issues[4][1] & High density, short\\
{}issues[4][2] & Kc mid < Kc ini\\
{}issues[5][0] & citrus\\
{}issues[5][1] & High density, tall\\
{}issues[5][2] & Kc mid < Kc ini\\
{}issues[6][0] & olive\\
{}issues[6][1] & Young\\
{}issues[6][2] & Kc mid < Kc ini, Kc mid 0.25\\
{}issues[7][0] & olive\\
{}issues[7][1] & Traditional, low density\\
{}issues[7][2] & Kc mid < Kc ini, Kc mid 0.3\\
{}issues[8][0] & olive\\
{}issues[8][1] & Traditional, medium density\\
{}issues[8][2] & Kc mid < Kc ini\\
{}issues[9][0] & olive\\
{}issues[9][1] & Intensive (hedge pruned)\\
{}issues[9][2] & Kc mid < Kc ini\\
{}issues[10][0] & olive\\
{}issues[10][1] & Super-intensive, medium\\
{}issues[10][2] & Kc mid < Kc ini\\
{}issues[11][0] & olive\\
{}issues[11][1] & Super-intensive, high\\
{}issues[11][2] & Kc mid < Kc ini\\
{}issues[12][0] & datepalm\\
{}issues[12][1] & Young, low\\
{}issues[12][2] & Kc mid < Kc ini\\
{}rows[0].key & rice\\
{}rows[0].dens & None\\
{}rows[0].group & field\\
{}rows[0].habit & annual\\
{}rows[0].kcIni & 1.05\\
{}rows[0].kcMid & 1.2\\
{}rows[0].kcEnd & 0.9\\
{}rows[0].len & 120\\
{}rows[0].root & 75\\
{}rows[0].p & 0.2\\
{}rows[0].ky & None\\
{}rows[0].kcSum & 135.1\\
{}rows[0].iss & \\
{}rows[1].key & maize\\
{}rows[1].dens & None\\
{}rows[1].group & field\\
{}rows[1].habit & annual\\
{}rows[1].kcIni & 0.3\\
{}rows[1].kcMid & 1.2\\
{}rows[1].kcEnd & 0.6\\
{}rows[1].len & 105\\
{}rows[1].root & 135\\
{}rows[1].p & 0.55\\
{}rows[1].ky & 1.25\\
{}rows[1].kcSum & 92.85\\
{}rows[1].iss & \\
{}rows[2].key & pearlmillet\\
{}rows[2].dens & None\\
{}rows[2].group & field\\
{}rows[2].habit & annual\\
{}rows[2].kcIni & 0.3\\
{}rows[2].kcMid & 1\\
{}rows[2].kcEnd & 0.3\\
{}rows[2].len & 105\\
{}rows[2].root & 150\\
{}rows[2].p & 0.55\\
{}rows[2].ky & None\\
{}rows[2].kcSum & 77\\
{}rows[2].iss & \\
{}rows[3].key & sorghum\\
{}rows[3].dens & None\\
{}rows[3].group & field\\
{}rows[3].habit & annual\\
{}rows[3].kcIni & 0.3\\
{}rows[3].kcMid & 1\\
{}rows[3].kcEnd & 0.55\\
{}rows[3].len & 125\\
{}rows[3].root & 150\\
{}rows[3].p & 0.55\\
{}rows[3].ky & 0.9\\
{}rows[3].kcSum & 91.88\\
{}rows[3].iss & \\
{}rows[4].key & cassava\\
{}rows[4].dens & None\\
{}rows[4].group & field\\
{}rows[4].habit & annual\\
{}rows[4].kcIni & 0.3\\
{}rows[4].kcMid & 0.8\\
{}rows[4].kcEnd & 0.3\\
{}rows[4].len & 210\\
{}rows[4].root & 65\\
{}rows[4].p & 0.35\\
{}rows[4].ky & None\\
{}rows[4].kcSum & 133\\
{}rows[4].iss & \\
{}rows[5].key & yam\\
{}rows[5].dens & None\\
{}rows[5].group & field\\
{}rows[5].habit & annual\\
{}rows[5].kcIni & 0.5\\
{}rows[5].kcMid & 1.1\\
{}rows[5].kcEnd & 0.85\\
{}rows[5].len & 240\\
{}rows[5].root & 60\\
{}rows[5].p & 0.4\\
{}rows[5].ky & None\\
{}rows[5].kcSum & 220.3\\
{}rows[5].iss & \\
{}rows[6].key & cowpea\\
{}rows[6].dens & None\\
{}rows[6].group & field\\
{}rows[6].habit & annual\\
{}rows[6].kcIni & 0.4\\
{}rows[6].kcMid & 1.05\\
{}rows[6].kcEnd & 0.5\\
{}rows[6].len & 100\\
{}rows[6].root & 80\\
{}rows[6].p & 0.45\\
{}rows[6].ky & None\\
{}rows[6].kcSum & 76.7\\
{}rows[6].iss & \\
{}rows[7].key & greengram\\
{}rows[7].dens & None\\
{}rows[7].group & field\\
{}rows[7].habit & annual\\
{}rows[7].kcIni & 0.4\\
{}rows[7].kcMid & 0.96\\
{}rows[7].kcEnd & 0.51\\
{}rows[7].len & 90\\
{}rows[7].root & 75\\
{}rows[7].p & 0.45\\
{}rows[7].ky & None\\
{}rows[7].kcSum & 66.78\\
{}rows[7].iss & \\
{}rows[8].key & soybean\\
{}rows[8].dens & None\\
{}rows[8].group & field\\
{}rows[8].habit & annual\\
{}rows[8].kcIni & 0.4\\
{}rows[8].kcMid & 1.15\\
{}rows[8].kcEnd & 0.5\\
{}rows[8].len & 120\\
{}rows[8].root & 95\\
{}rows[8].p & 0.5\\
{}rows[8].ky & 0.85\\
{}rows[8].kcSum & 103.6\\
{}rows[8].iss & \\
{}rows[9].key & groundnut\\
{}rows[9].dens & None\\
{}rows[9].group & field\\
{}rows[9].habit & annual\\
{}rows[9].kcIni & 0.4\\
{}rows[9].kcMid & 1.15\\
{}rows[9].kcEnd & 0.6\\
{}rows[9].len & 120\\
{}rows[9].root & 75\\
{}rows[9].p & 0.5\\
{}rows[9].ky & 0.7\\
{}rows[9].kcSum & 102.9\\
{}rows[9].iss & \\
{}rows[10].key & cotton\\
{}rows[10].dens & None\\
{}rows[10].group & field\\
{}rows[10].habit & annual\\
{}rows[10].kcIni & 0.35\\
{}rows[10].kcMid & 1.15\\
{}rows[10].kcEnd & 0.5\\
{}rows[10].len & 190\\
{}rows[10].root & 135\\
{}rows[10].p & 0.65\\
{}rows[10].ky & 0.85\\
{}rows[10].kcSum & 160.2\\
{}rows[10].iss & \\
{}rows[11].key & wheat\\
{}rows[11].dens & None\\
{}rows[11].group & field\\
{}rows[11].habit & annual\\
{}rows[11].kcIni & 0.4\\
{}rows[11].kcMid & 1.15\\
{}rows[11].kcEnd & 0.41\\
{}rows[11].len & 120\\
{}rows[11].root & 165\\
{}rows[11].p & 0.55\\
{}rows[11].ky & 1.05\\
{}rows[11].kcSum & 102.5\\
{}rows[11].iss & \\
{}rows[12].key & chickpea\\
{}rows[12].dens & None\\
{}rows[12].group & field\\
{}rows[12].habit & annual\\
{}rows[12].kcIni & 0.4\\
{}rows[12].kcMid & 1\\
{}rows[12].kcEnd & 0.4\\
{}rows[12].len & 120\\
{}rows[12].root & 80\\
{}rows[12].p & 0.5\\
{}rows[12].ky & None\\
{}rows[12].kcSum & 93\\
{}rows[12].iss & \\
{}rows[13].key & mustard\\
{}rows[13].dens & None\\
{}rows[13].group & field\\
{}rows[13].habit & annual\\
{}rows[13].kcIni & 0.35\\
\bottomrule\end{longtable}}

\subsection*{out\_p6\_pheno.json}
\noindent\small Source: \nolinkurl{Manuscripts/reproducibility/p456/out_p6_pheno.json}; 560 values, first 220 shown.\par
{\scriptsize\begin{longtable}{p{\dimexpr 0.4557\linewidth-2\tabcolsep\relax}p{\dimexpr 0.5443\linewidth-2\tabcolsep\relax}}\toprule\textbf{quantity} & \textbf{value}\\\midrule\endhead
{}[0].clim & warm\\
{}[0].geno & Reference\\
{}[0].tr & ww\\
{}[0].closure\_mm & -6.724e-07\\
{}[0].det & 0\\
{}[0].P & 0\\
{}[0].I & 589.1\\
{}[0].T & 330\\
{}[0].Es & 252.3\\
{}[0].drain & 0\\
{}[0].grain & 515.1\\
{}[0].shoot & 542.3\\
{}[0].root & 172.6\\
{}[0].rdepth & 119.5\\
{}[0].cwsiMax & 0.2467\\
{}[0].ctdMax & 1.066\\
{}[0].psiMin & -1.213\\
{}[0].lai & 2.632\\
{}[0].height & 213.7\\
{}[0].stage & 3.657\\
{}[1].clim & warm\\
{}[1].geno & Reference\\
{}[1].tr & dr\\
{}[1].closure\_mm & 2.013e-05\\
{}[1].det & 0\\
{}[1].P & 0\\
{}[1].I & 519.7\\
{}[1].T & 253.9\\
{}[1].Es & 259.1\\
{}[1].drain & 1.907e-06\\
{}[1].grain & 345.4\\
{}[1].shoot & 478.6\\
{}[1].root & 166.1\\
{}[1].rdepth & 119.5\\
{}[1].cwsiMax & 0.3367\\
{}[1].ctdMax & 1.076\\
{}[1].psiMin & -1.285\\
{}[1].lai & 2.429\\
{}[1].height & 213\\
{}[1].stage & 3.657\\
{}[2].clim & warm\\
{}[2].geno & Deep rooting\\
{}[2].tr & ww\\
{}[2].closure\_mm & -2.057e-05\\
{}[2].det & 0\\
{}[2].P & 0\\
{}[2].I & 584.2\\
{}[2].T & 305.8\\
{}[2].Es & 271.6\\
{}[2].drain & 0\\
{}[2].grain & 485.8\\
{}[2].shoot & 481.6\\
{}[2].root & 186.1\\
{}[2].rdepth & 119.5\\
{}[2].cwsiMax & 0.2172\\
{}[2].ctdMax & 1.078\\
{}[2].psiMin & -1.203\\
{}[2].lai & 2.326\\
{}[2].height & 213.7\\
{}[2].stage & 3.657\\
{}[3].clim & warm\\
{}[3].geno & Deep rooting\\
{}[3].tr & dr\\
{}[3].closure\_mm & -3.866e-05\\
{}[3].det & 0\\
{}[3].P & 0\\
{}[3].I & 518.4\\
{}[3].T & 239.2\\
{}[3].Es & 272.4\\
{}[3].drain & 0\\
{}[3].grain & 328.9\\
{}[3].shoot & 432.8\\
{}[3].root & 179.2\\
{}[3].rdepth & 119.5\\
{}[3].cwsiMax & 0.346\\
{}[3].ctdMax & 1.085\\
{}[3].psiMin & -1.282\\
{}[3].lai & 2.192\\
{}[3].height & 213.1\\
{}[3].stage & 3.657\\
{}[4].clim & warm\\
{}[4].geno & Shallow, dense\\
{}[4].tr & ww\\
{}[4].closure\_mm & -1.514e-05\\
{}[4].det & 0\\
{}[4].P & 0\\
{}[4].I & 589.6\\
{}[4].T & 331.3\\
{}[4].Es & 251.6\\
{}[4].drain & 0\\
{}[4].grain & 516.2\\
{}[4].shoot & 544\\
{}[4].root & 172.7\\
{}[4].rdepth & 109\\
{}[4].cwsiMax & 0.2477\\
{}[4].ctdMax & 1.066\\
{}[4].psiMin & -1.214\\
{}[4].lai & 2.644\\
{}[4].height & 213.6\\
{}[4].stage & 3.657\\
{}[5].clim & warm\\
{}[5].geno & Shallow, dense\\
{}[5].tr & dr\\
{}[5].closure\_mm & 2.381e-05\\
{}[5].det & 0\\
{}[5].P & 0\\
{}[5].I & 501.3\\
{}[5].T & 228.3\\
{}[5].Es & 266.3\\
{}[5].drain & 1.907e-06\\
{}[5].grain & 297.7\\
{}[5].shoot & 429.1\\
{}[5].root & 158.4\\
{}[5].rdepth & 111.5\\
{}[5].cwsiMax & 0.3839\\
{}[5].ctdMax & 1.082\\
{}[5].psiMin & -1.27\\
{}[5].lai & 2.219\\
{}[5].height & 212.7\\
{}[5].stage & 3.657\\
{}[6].clim & warm\\
{}[6].geno & Stay-green\\
{}[6].tr & ww\\
{}[6].closure\_mm & 1.436e-05\\
{}[6].det & 0\\
{}[6].P & 0\\
{}[6].I & 611.2\\
{}[6].T & 395\\
{}[6].Es & 209.6\\
{}[6].drain & 0\\
{}[6].grain & 541.7\\
{}[6].shoot & 552.9\\
{}[6].root & 176.6\\
{}[6].rdepth & 119.5\\
{}[6].cwsiMax & 0.2492\\
{}[6].ctdMax & 1.113\\
{}[6].psiMin & -1.211\\
{}[6].lai & 2.738\\
{}[6].height & 213.7\\
{}[6].stage & 3.657\\
{}[7].clim & warm\\
{}[7].geno & Stay-green\\
{}[7].tr & dr\\
{}[7].closure\_mm & 4.894e-06\\
{}[7].det & 0\\
{}[7].P & 0\\
{}[7].I & 543.4\\
{}[7].T & 319\\
{}[7].Es & 217.8\\
{}[7].drain & 0\\
{}[7].grain & 376.9\\
{}[7].shoot & 488.7\\
{}[7].root & 170.7\\
{}[7].rdepth & 119.5\\
{}[7].cwsiMax & 0.3559\\
{}[7].ctdMax & 1.11\\
{}[7].psiMin & -1.265\\
{}[7].lai & 2.576\\
{}[7].height & 213\\
{}[7].stage & 3.657\\
{}[8].clim & warm\\
{}[8].geno & Water saver\\
{}[8].tr & ww\\
{}[8].closure\_mm & -1.134e-05\\
{}[8].det & 0\\
{}[8].P & 0\\
{}[8].I & 506.2\\
{}[8].T & 164.5\\
{}[8].Es & 334.8\\
{}[8].drain & 0\\
{}[8].grain & 330.9\\
{}[8].shoot & 336\\
{}[8].root & 143.3\\
{}[8].rdepth & 119.5\\
{}[8].cwsiMax & 0.3778\\
{}[8].ctdMax & 0.9156\\
{}[8].psiMin & -0.8846\\
{}[8].lai & 1.48\\
{}[8].height & 213.7\\
{}[8].stage & 3.657\\
{}[9].clim & warm\\
{}[9].geno & Water saver\\
{}[9].tr & dr\\
{}[9].closure\_mm & -6.134e-06\\
{}[9].det & 0\\
{}[9].P & 0\\
{}[9].I & 453.8\\
{}[9].T & 135\\
{}[9].Es & 311.9\\
{}[9].drain & 0\\
{}[9].grain & 241.9\\
{}[9].shoot & 300.1\\
{}[9].root & 135.6\\
{}[9].rdepth & 119.5\\
{}[9].cwsiMax & 0.4086\\
{}[9].ctdMax & 0.9363\\
{}[9].psiMin & -0.9405\\
{}[9].lai & 1.364\\
{}[9].height & 213.2\\
{}[9].stage & 3.657\\
{}[10].clim & warm\\
{}[10].geno & Water spender\\
{}[10].tr & ww\\
{}[10].closure\_mm & 1.024e-05\\
{}[10].det & 0\\
{}[10].P & 0\\
{}[10].I & 615.1\\
{}[10].T & 377\\
{}[10].Es & 231.5\\
{}[10].drain & 0\\
{}[10].grain & 561.5\\
{}[10].shoot & 608\\
{}[10].root & 180.6\\
{}[10].rdepth & 119.5\\
{}[10].cwsiMax & 0.1435\\
{}[10].ctdMax & 1.114\\
{}[10].psiMin & -1.59\\
{}[10].lai & 3.01\\
{}[10].height & 213.5\\
{}[10].stage & 3.657\\
\bottomrule\end{longtable}}

\subsection*{out\_p6\_sweep.json}
\noindent\small Source: \nolinkurl{Manuscripts/reproducibility/p456/out_p6_sweep.json}; 286 values, first 220 shown.\par
{\scriptsize\begin{longtable}{p{\dimexpr 0.4557\linewidth-2\tabcolsep\relax}p{\dimexpr 0.5443\linewidth-2\tabcolsep\relax}}\toprule\textbf{quantity} & \textbf{value}\\\midrule\endhead
{}[0].id & sid\\
{}[0].ran & False\\
{}[0].engineInputs & 19\\
{}[0].ms & 6655\\
{}[0].tabs & Results, Design, Charts, Figure, Messages\\
{}[0].inputs & 0\\
{}[0].chartsTab & 8\\
{}[0].errors & \\
{}[0].kpi & 62\\
{}[0].msgs & 0\\
{}[0].charts & 4\\
{}[0].textLen & 2259\\
{}[0].canvas & None\\
{}[1].id & drip\\
{}[1].ran & False\\
{}[1].engineInputs & 19\\
{}[1].ms & 6639\\
{}[1].tabs & Results, Design, Charts, Figure, Messages\\
{}[1].inputs & 0\\
{}[1].chartsTab & 8\\
{}[1].errors & \\
{}[1].kpi & 62\\
{}[1].msgs & 2\\
{}[1].charts & 4\\
{}[1].textLen & 2500\\
{}[1].canvas & None\\
{}[2].id & spr\\
{}[2].ran & False\\
{}[2].engineInputs & 19\\
{}[2].ms & 6839\\
{}[2].tabs & Results, Design, Charts, Figure, Messages\\
{}[2].inputs & 0\\
{}[2].chartsTab & 8\\
{}[2].errors & \\
{}[2].kpi & 62\\
{}[2].msgs & 4\\
{}[2].charts & 4\\
{}[2].textLen & 2860\\
{}[2].canvas & None\\
{}[3].id & v1\\
{}[3].ran & True\\
{}[3].engineInputs & 19\\
{}[3].ms & 6431\\
{}[3].tabs & Results, Design, Charts, Figure, Messages\\
{}[3].inputs & 0\\
{}[3].chartsTab & 4\\
{}[3].errors & \\
{}[3].kpi & 62\\
{}[3].msgs & 4\\
{}[3].charts & 4\\
{}[3].textLen & 1672\\
{}[3].canvas & None\\
{}[4].id & n1\\
{}[4].ran & True\\
{}[4].engineInputs & 19\\
{}[4].ms & 6350\\
{}[4].tabs & Results, Design, Charts, Figure, Messages\\
{}[4].inputs & 0\\
{}[4].chartsTab & 4\\
{}[4].errors & \\
{}[4].kpi & 62\\
{}[4].msgs & 6\\
{}[4].charts & 6\\
{}[4].textLen & 3419\\
{}[4].canvas & 606, 460\\
{}[5].id & v4\\
{}[5].ran & True\\
{}[5].engineInputs & 19\\
{}[5].ms & 6388\\
{}[5].tabs & Results, Design, Charts, Figure, Messages\\
{}[5].inputs & 0\\
{}[5].chartsTab & 4\\
{}[5].errors & \\
{}[5].kpi & 62\\
{}[5].msgs & 4\\
{}[5].charts & 4\\
{}[5].textLen & 1584\\
{}[5].canvas & None\\
{}[6].id & v6\\
{}[6].ran & True\\
{}[6].engineInputs & 19\\
{}[6].ms & 6371\\
{}[6].tabs & Results, Design, Charts, Figure, Messages\\
{}[6].inputs & 0\\
{}[6].chartsTab & 4\\
{}[6].errors & \\
{}[6].kpi & 62\\
{}[6].msgs & 2\\
{}[6].charts & 4\\
{}[6].textLen & 1503\\
{}[6].canvas & None\\
{}[7].id & e1\\
{}[7].ran & True\\
{}[7].engineInputs & 19\\
{}[7].ms & 6574\\
{}[7].tabs & Results, Design, Charts, Figure, Messages\\
{}[7].inputs & 0\\
{}[7].chartsTab & 4\\
{}[7].errors & \\
{}[7].kpi & 62\\
{}[7].msgs & 3\\
{}[7].charts & 4\\
{}[7].textLen & 1802\\
{}[7].canvas & None\\
{}[8].id & n2\\
{}[8].ran & True\\
{}[8].engineInputs & 19\\
{}[8].ms & 6405\\
{}[8].tabs & Results, Design, Charts, Figure, Messages\\
{}[8].inputs & 0\\
{}[8].chartsTab & 4\\
{}[8].errors & \\
{}[8].kpi & 62\\
{}[8].msgs & 2\\
{}[8].charts & 4\\
{}[8].textLen & 1674\\
{}[8].canvas & None\\
{}[9].id & n3\\
{}[9].ran & True\\
{}[9].engineInputs & 19\\
{}[9].ms & 6370\\
{}[9].tabs & Results, Design, Charts, Figure, Messages\\
{}[9].inputs & 0\\
{}[9].chartsTab & 4\\
{}[9].errors & \\
{}[9].kpi & 62\\
{}[9].msgs & 5\\
{}[9].charts & 4\\
{}[9].textLen & 1680\\
{}[9].canvas & None\\
{}[10].id & v3\\
{}[10].ran & True\\
{}[10].engineInputs & 19\\
{}[10].ms & 6457\\
{}[10].tabs & Results, Design, Charts, Figure, Messages\\
{}[10].inputs & 0\\
{}[10].chartsTab & 4\\
{}[10].errors & \\
{}[10].kpi & 62\\
{}[10].msgs & 3\\
{}[10].charts & 4\\
{}[10].textLen & 1661\\
{}[10].canvas & None\\
{}[11].id & v7\\
{}[11].ran & True\\
{}[11].engineInputs & 19\\
{}[11].ms & 6461\\
{}[11].tabs & Results, Design, Charts, Figure, Messages\\
{}[11].inputs & 0\\
{}[11].chartsTab & 4\\
{}[11].errors & \\
{}[11].kpi & 62\\
{}[11].msgs & 0\\
{}[11].charts & 6\\
{}[11].textLen & 3237\\
{}[11].canvas & 606, 440\\
{}[12].id & r1\\
{}[12].ran & True\\
{}[12].engineInputs & 19\\
{}[12].ms & 6429\\
{}[12].tabs & Results, Design, Charts, Figure, Messages\\
{}[12].inputs & 0\\
{}[12].chartsTab & 4\\
{}[12].errors & \\
{}[12].kpi & 62\\
{}[12].msgs & 3\\
{}[12].charts & 4\\
{}[12].textLen & 1987\\
{}[12].canvas & None\\
{}[13].id & dph\\
{}[13].ran & False\\
{}[13].engineInputs & 19\\
{}[13].ms & 6328\\
{}[13].tabs & Results, Design, Charts, Figure, Messages\\
{}[13].inputs & 0\\
{}[13].chartsTab & 0\\
{}[13].errors & \\
{}[13].kpi & 0\\
{}[13].msgs & 0\\
{}[13].charts & 0\\
{}[13].textLen & 292\\
{}[13].canvas & None\\
{}[14].id & field\\
{}[14].ran & True\\
{}[14].engineInputs & 19\\
{}[14].ms & 6716\\
{}[14].tabs & Results, Design, Charts, Figure, Messages\\
{}[14].inputs & 0\\
{}[14].chartsTab & 4\\
{}[14].errors & \\
{}[14].kpi & 62\\
{}[14].msgs & 5\\
{}[14].charts & 4\\
{}[14].textLen & 2084\\
{}[14].canvas & None\\
{}[15].id & spac\\
{}[15].ran & True\\
{}[15].engineInputs & 19\\
{}[15].ms & 6562\\
{}[15].tabs & Results, Design, Charts, Figure, Messages\\
{}[15].inputs & 0\\
{}[15].chartsTab & 4\\
{}[15].errors & \\
{}[15].kpi & 62\\
{}[15].msgs & 0\\
{}[15].charts & 4\\
{}[15].textLen & 1470\\
{}[15].canvas & None\\
{}[16].id & v2\\
{}[16].ran & True\\
{}[16].engineInputs & 19\\
{}[16].ms & 6490\\
{}[16].tabs & Results, Design, Charts, Figure, Messages\\
{}[16].inputs & 0\\
{}[16].chartsTab & 4\\
{}[16].errors & \\
{}[16].kpi & 62\\
{}[16].msgs & 4\\
{}[16].charts & 4\\
{}[16].textLen & 1695\\
\bottomrule\end{longtable}}

\subsection*{out\_p7\_planner.json}
\noindent\small Source: \nolinkurl{Manuscripts/reproducibility/p456/out_p7_planner.json}; 57 values.\par
{\scriptsize\begin{longtable}{p{\dimexpr 0.4557\linewidth-2\tabcolsep\relax}p{\dimexpr 0.5443\linewidth-2\tabcolsep\relax}}\toprule\textbf{quantity} & \textbf{value}\\\midrule\endhead
{}n\_questions & 8\\
{}passed & 8\\
{}rows[0].q & How does the irrigation trigger of -100 to -800 cm change drainage on a sandy loam in Bologna over 10 years?\\
{}rows[0].expected & trig\\
{}rows[0].got & irrig, trig\\
{}rows[0].station & it-bologna\\
{}rows[0].soil & usda:sandyloam\\
{}rows[0].pass & True\\
{}rows[1].q & What happens to nitrate leaching if I cut the nitrogen dose by 20 \% in Foulum?\\
{}rows[1].expected & nfac\\
{}rows[1].got & nfac, mm\\
{}rows[1].station & dk-foulum\\
{}rows[1].soil & None\\
{}rows[1].pass & True\\
{}rows[2].q & Compare 20 \% less rain and 20 \% more rain at Jyndevad on a loamy sand\\
{}rows[2].expected & precip\\
{}rows[2].got & precip\\
{}rows[2].station & dk-jyndevad\\
{}rows[2].soil & usda:loamysand\\
{}rows[2].pass & True\\
{}rows[3].q & Warming of 2 degrees and 4 degrees in Poland, Jadwisin\\
{}rows[3].expected & temp\\
{}rows[3].got & temp\\
{}rows[3].station & pl-jadwisin\\
{}rows[3].soil & None\\
{}rows[3].pass & True\\
{}rows[4].q & Is drip irrigation better than rainfed on a silt loam in the Czech Republic?\\
{}rows[4].expected & irrig\\
{}rows[4].got & irrig, method\\
{}rows[4].station & cz-valecov\\
{}rows[4].soil & usda:siltloam\\
{}rows[4].pass & True\\
{}rows[5].q & Effect of the sensor depth 10, 30 and 60 cm on the irrigation trigger\\
{}rows[5].expected & z\\
{}rows[5].got & irrig, trig, z\\
{}rows[5].station & None\\
{}rows[5].soil & None\\
{}rows[5].pass & True\\
{}rows[6].q & Pesticide half-life of 20 and 60 days and sorption koc of 50 and 200 l/kg\\
{}rows[6].expected & half, koc\\
{}rows[6].got & koc, half\\
{}rows[6].station & None\\
{}rows[6].soil & None\\
{}rows[6].pass & True\\
{}rows[7].q & What if the rooting depth is 60 and 120 cm in Bratislava?\\
{}rows[7].expected & root\\
{}rows[7].got & root\\
{}rows[7].station & sk-bratislava\\
{}rows[7].soil & None\\
{}rows[7].pass & True\\
{}catalogue\_factors & 54\\
{}catalogue\_soils & 10\\
{}catalogue\_stations & 8\\
{}validate.factors[0].id & precip\\
{}validate.factors[0].levels & 0.4, 1, 1.6\\
{}validate.responses & yield\\
{}validate.notes & unknown factor "not/\_a/\_factor" left out\\
\bottomrule\end{longtable}}


\section{Where every table and figure comes from}
{\scriptsize\begin{longtable}{p{\dimexpr 0.0955\linewidth-2\tabcolsep\relax}p{\dimexpr 0.2293\linewidth-2\tabcolsep\relax}p{\dimexpr 0.3185\linewidth-2\tabcolsep\relax}p{\dimexpr 0.3567\linewidth-2\tabcolsep\relax}}\toprule\textbf{Item} & \textbf{How it is made} & \textbf{Script(s)} & \textbf{Input and output}\\\midrule\endhead
{}Figure 1 & Figure drawn by the script from the input data & Manuscripts/\allowbreak{}reproducibility/\allowbreak{}p456/\allowbreak{}p456\_\allowbreak{}figs\_\allowbreak{}architecture.py & Data/\allowbreak{}daisy\_\allowbreak{}ensemble.csv.gz (10 190 runs);  Manuscripts/\allowbreak{}figures/\allowbreak{}p6/\allowbreak{}fig01\_\allowbreak{}system.pdf\\
{}Table 1 & Descriptive table (definitions, design, parameters or tests); no computation & - & \\
{}Figure 2 & Screen capture of SWAFL Desktop 1.1.0 (not regenerated) & - & Manuscripts/\allowbreak{}figures/\allowbreak{}shots\_\allowbreak{}jpg/\allowbreak{}zaria\_\allowbreak{}ndvi.jpg;  Manuscripts/\allowbreak{}figures/\allowbreak{}shots\_\allowbreak{}jpg/\allowbreak{}zaria\_\allowbreak{}scenes.jpg;  Manuscripts/\allowbreak{}figures/\allowbreak{}shots\_\allowbreak{}jpg/\allowbreak{}real\_\allowbreak{}sketch.jpg\\
{}Figure 3 & Screen capture of SWAFL Desktop 1.1.0 (not regenerated) & - & Manuscripts/\allowbreak{}figures/\allowbreak{}shots\_\allowbreak{}jpg/\allowbreak{}zaria\_\allowbreak{}terrain.jpg;  Manuscripts/\allowbreak{}figures/\allowbreak{}shots\_\allowbreak{}jpg/\allowbreak{}zaria\_\allowbreak{}gndvi\_\allowbreak{}water.jpg;  Manuscripts/\allowbreak{}figures/\allowbreak{}shots\_\allowbreak{}jpg/\allowbreak{}zaria\_\allowbreak{}rusle.jpg;  Manuscripts/\allowbreak{}figures/\allowbreak{}shots\_\allowbreak{}jpg/\allowbreak{}zaria\_\allowbreak{}available\_\allowbreak{}water.jpg\\
{}Figure 4 & Screen capture of SWAFL Desktop 1.1.0 (not regenerated) & - & Manuscripts/\allowbreak{}figures/\allowbreak{}shots\_\allowbreak{}jpg/\allowbreak{}zaria\_\allowbreak{}vri\_\allowbreak{}depth.jpg\\
{}Figure 5 & Screen capture of SWAFL Desktop 1.1.0 (not regenerated) & - & Manuscripts/\allowbreak{}figures/\allowbreak{}p6/\allowbreak{}fig03\_\allowbreak{}pheno\_\allowbreak{}figures.jpg\\
{}Figure 6 & Screen capture of SWAFL Desktop 1.1.0 (not regenerated) & - & Manuscripts/\allowbreak{}figures/\allowbreak{}shots\_\allowbreak{}jpg/\allowbreak{}real\_\allowbreak{}season\_\allowbreak{}film.jpg\\
{}Figure 7 & Figure drawn by the script from the input data & Manuscripts/\allowbreak{}reproducibility/\allowbreak{}p456/\allowbreak{}p456\_\allowbreak{}figs\_\allowbreak{}architecture.py & Data/\allowbreak{}daisy\_\allowbreak{}ensemble.csv.gz (10 190 runs);  Manuscripts/\allowbreak{}figures/\allowbreak{}p4/\allowbreak{}fig01\_\allowbreak{}workbench.pdf\\
{}Figure 8 & Screen capture of SWAFL Desktop 1.1.0 (not regenerated) & - & Manuscripts/\allowbreak{}figures/\allowbreak{}p4/\allowbreak{}fig02a\_\allowbreak{}process\_\allowbreak{}map.png\\
{}Figure 9 & Screen capture of SWAFL Desktop 1.1.0 (not regenerated) & - & Manuscripts/\allowbreak{}figures/\allowbreak{}p4/\allowbreak{}fig02b\_\allowbreak{}soil\_\allowbreak{}column.png\\
{}Figure 10 & Screen capture of SWAFL Desktop 1.1.0 (not regenerated) & - & Manuscripts/\allowbreak{}figures/\allowbreak{}shots\_\allowbreak{}jpg/\allowbreak{}wb\_\allowbreak{}step1\_\allowbreak{}input.jpg;  Manuscripts/\allowbreak{}figures/\allowbreak{}shots\_\allowbreak{}jpg/\allowbreak{}wb\_\allowbreak{}step2\_\allowbreak{}library.jpg;  Manuscripts/\allowbreak{}figures/\allowbreak{}shots\_\allowbreak{}jpg/\allowbreak{}wb\_\allowbreak{}step4\_\allowbreak{}scenarios.jpg;  Manuscripts/\allowbreak{}figures/\allowbreak{}shots\_\allowbreak{}jpg/\allowbreak{}wb\_\allowbreak{}step6\_\allowbreak{}logplots.jpg;  Manuscripts/\allowbreak{}figures/\allowbreak{}shots\_\allowbreak{}jpg/\allowbreak{}wb\_\allowbreak{}step9\_\allowbreak{}layers.jpg;  Manuscripts/\allowbreak{}figures/\allowbreak{}shots\_\allowbreak{}jpg/\allowbreak{}wb\_\allowbreak{}tool\_\allowbreak{}setup\_\allowbreak{}editor.jpg;  Manuscripts/\allowbreak{}figures/\allowbreak{}shots\_\allowbreak{}jpg/\allowbreak{}wb\_\allowbreak{}tool\_\allowbreak{}calibrate.jpg;  Manuscripts/\allowbreak{}figures/\allowbreak{}shots\_\allowbreak{}jpg/\allowbreak{}wb\_\allowbreak{}tool\_\allowbreak{}learn\_\allowbreak{}files.jpg\\
{}Figure 11 & Screen capture of SWAFL Desktop 1.1.0 (not regenerated) & - & Manuscripts/\allowbreak{}figures/\allowbreak{}shots\_\allowbreak{}jpg/\allowbreak{}zaria\_\allowbreak{}main\_\allowbreak{}feedback.jpg\\
{}Figure 12 & Figure drawn by the script from the input data & Manuscripts/\allowbreak{}reproducibility/\allowbreak{}p456/\allowbreak{}p456\_\allowbreak{}figs\_\allowbreak{}schematics.py & Data/\allowbreak{}daisy\_\allowbreak{}ensemble.csv.gz (10 190 runs);  Manuscripts/\allowbreak{}figures/\allowbreak{}p5/\allowbreak{}fig01\_\allowbreak{}bridge.pdf\\
{}Table 2 & Descriptive table (definitions, design, parameters or tests); no computation & - & \\
{}Table 3 & Table typed from the result file (13 of its numbers found there) & Manuscripts/\allowbreak{}reproducibility/\allowbreak{}p3\_\allowbreak{}js/\allowbreak{}run\_\allowbreak{}drip.js;  Manuscripts/\allowbreak{}reproducibility/\allowbreak{}analysis/\allowbreak{}p3\_\allowbreak{}lit\_\allowbreak{}figs.py;  Manuscripts/\allowbreak{}reproducibility/\allowbreak{}analysis/\allowbreak{}make\_\allowbreak{}supplement\_\allowbreak{}data.py & Data/\allowbreak{}daisy\_\allowbreak{}ensemble.csv.gz (10 190 runs);  Manuscripts/\allowbreak{}results/\allowbreak{}drip\_\allowbreak{}out.json\\
{}Table 4 & Table typed in the manuscript from the result files of the script & p4\_\allowbreak{}design\_\allowbreak{}sizes.js;  out\_\allowbreak{}p4\_\allowbreak{}design.js & Data/\allowbreak{}daisy\_\allowbreak{}ensemble.csv.gz (10 190 runs)\\
{}Figure 13 & Figure drawn by the script from the input data & Manuscripts/\allowbreak{}reproducibility/\allowbreak{}p456/\allowbreak{}p456\_\allowbreak{}figs\_\allowbreak{}data.py;  p4\_\allowbreak{}design\_\allowbreak{}fidelity.py & Data/\allowbreak{}daisy\_\allowbreak{}ensemble.csv.gz (10 190 runs);  Manuscripts/\allowbreak{}figures/\allowbreak{}p4/\allowbreak{}fig03\_\allowbreak{}fidelity.pdf\\
{}Table 5 & Table typed from the result file (5 of its numbers found there) & Manuscripts/\allowbreak{}reproducibility/\allowbreak{}run\_\allowbreak{}all.sh;  Manuscripts/\allowbreak{}reproducibility/\allowbreak{}make\_\allowbreak{}e12.py;  Manuscripts/\allowbreak{}reproducibility/\allowbreak{}audit\_\allowbreak{}numbers.py;  p4\_\allowbreak{}richards.js;  p4\_\allowbreak{}richards\_\allowbreak{}ref.py & Data/\allowbreak{}daisy\_\allowbreak{}ensemble.csv.gz (10 190 runs);  Manuscripts/\allowbreak{}results/\allowbreak{}e12.json\\
{}Figure 14 & Figure drawn by the script from the input data & Manuscripts/\allowbreak{}reproducibility/\allowbreak{}p456/\allowbreak{}p456\_\allowbreak{}figs\_\allowbreak{}data.py;  p4\_\allowbreak{}richards.js;  p4\_\allowbreak{}richards\_\allowbreak{}ref.py & Data/\allowbreak{}daisy\_\allowbreak{}ensemble.csv.gz (10 190 runs);  Manuscripts/\allowbreak{}figures/\allowbreak{}p4/\allowbreak{}fig04\_\allowbreak{}richards.pdf\\
{}Figure 15 & Figure drawn by the script from the input data & Manuscripts/\allowbreak{}reproducibility/\allowbreak{}p456/\allowbreak{}p456\_\allowbreak{}figs\_\allowbreak{}data.py;  p5\_\allowbreak{}bridge\_\allowbreak{}checks.js & Data/\allowbreak{}daisy\_\allowbreak{}ensemble.csv.gz (10 190 runs);  Manuscripts/\allowbreak{}figures/\allowbreak{}p4/\allowbreak{}fig05\_\allowbreak{}denitrification.pdf\\
{}Table 6 & Descriptive table (definitions, design, parameters or tests); no computation & - & \\
{}Table 7 & Table typed in the manuscript from the result files of the script & (see the reproducibility appendix of the paper) & Data/\allowbreak{}daisy\_\allowbreak{}ensemble.csv.gz (10 190 runs)\\
{}Table 8 & Descriptive table (definitions, design, parameters or tests); no computation & - & \\
{}Figure 16 & Figure drawn by the script from the input data & Manuscripts/\allowbreak{}reproducibility/\allowbreak{}p456/\allowbreak{}p456\_\allowbreak{}figs\_\allowbreak{}data.py & Data/\allowbreak{}daisy\_\allowbreak{}ensemble.csv.gz (10 190 runs);  Manuscripts/\allowbreak{}figures/\allowbreak{}p5/\allowbreak{}fig04\_\allowbreak{}economics.pdf\\
{}Figure 17 & Screen capture of SWAFL Desktop 1.1.0 (not regenerated) & - & Manuscripts/\allowbreak{}figures/\allowbreak{}p5/\allowbreak{}fig02\_\allowbreak{}flowfigures.jpg\\
{}Figure 18 & Screen capture of SWAFL Desktop 1.1.0 (not regenerated) & - & Manuscripts/\allowbreak{}figures/\allowbreak{}shots\_\allowbreak{}jpg/\allowbreak{}ui\_\allowbreak{}layers\_\allowbreak{}tree.jpg;  Manuscripts/\allowbreak{}figures/\allowbreak{}shots\_\allowbreak{}jpg/\allowbreak{}ui\_\allowbreak{}field\_\allowbreak{}panel1.jpg;  Manuscripts/\allowbreak{}figures/\allowbreak{}shots\_\allowbreak{}jpg/\allowbreak{}ui\_\allowbreak{}field\_\allowbreak{}panel2.jpg\\
{}Figure 19 & Screen capture of SWAFL Desktop 1.1.0 (not regenerated) & - & Manuscripts/\allowbreak{}figures/\allowbreak{}shots\_\allowbreak{}jpg/\allowbreak{}wb\_\allowbreak{}summary\_\allowbreak{}effects.jpg;  Manuscripts/\allowbreak{}figures/\allowbreak{}shots\_\allowbreak{}jpg/\allowbreak{}wb\_\allowbreak{}every\_\allowbreak{}log\_\allowbreak{}lai.jpg\\
{}Figure 20 & Screen capture of SWAFL Desktop 1.1.0 (not regenerated) & - & Manuscripts/\allowbreak{}figures/\allowbreak{}shots\_\allowbreak{}jpg/\allowbreak{}wb\_\allowbreak{}why\_\allowbreak{}failed.jpg;  Manuscripts/\allowbreak{}figures/\allowbreak{}shots\_\allowbreak{}jpg/\allowbreak{}wb\_\allowbreak{}all\_\allowbreak{}outputs.jpg\\
{}Table 9 & Table typed in the manuscript from the result files of the script & (see the reproducibility appendix of the paper) & Data/\allowbreak{}daisy\_\allowbreak{}ensemble.csv.gz (10 190 runs)\\
{}Figure 21 & Figure drawn by the script from the input data & Manuscripts/\allowbreak{}reproducibility/\allowbreak{}p456/\allowbreak{}p456\_\allowbreak{}figs\_\allowbreak{}data.py & Data/\allowbreak{}daisy\_\allowbreak{}ensemble.csv.gz (10 190 runs);  Manuscripts/\allowbreak{}figures/\allowbreak{}p6/\allowbreak{}fig04\_\allowbreak{}pheno.pdf\\
{}Figure 22 & Figure drawn by the script from the input data & Manuscripts/\allowbreak{}reproducibility/\allowbreak{}p456/\allowbreak{}p456\_\allowbreak{}figs\_\allowbreak{}data.py & Data/\allowbreak{}daisy\_\allowbreak{}ensemble.csv.gz (10 190 runs);  Manuscripts/\allowbreak{}figures/\allowbreak{}p6/\allowbreak{}fig05\_\allowbreak{}crops.pdf\\
{}Table 10 & Descriptive table (definitions, design, parameters or tests); no computation & - & \\
\bottomrule\end{longtable}}


\bibliography{refs}
\end{document}
